\begin{textsample}{POS Dim 3 – 2017\_c6 – Score 2.00 – 2017\_c6/enade\_2017\_ci\_bio\_3400d20d88.txt}  \label{ex:f3_pos_192}
CAPÍTULO 6 ANÁLISE TÉCNICA DA PROVA Este capítulo tem por objetivo apresentar o desempenho dos estudantes concluintes de Ciências Biológicas ( Licenciatura/Bacharelado ) no Enade/2017.
Para isso, foram calculadas as estatísticas básicas da prova em seu todo ( seção 6.1.1 ), bem como as estatísticas relacionadas aos Componentes de Formação Geral ( seção 6.1.2 ) e de Conhecimento Específico da Área ( seção 6.1.3 ).
Dadas as suas características, foram analisadas, em separado, as questões objetivas ( seção 6.2 ) e as questões discursivas ( seção 6.3 ).
Para as questões objetivas, foram disponibilizados os índices de facilidade e de discriminação ponto-bisserial, também em separado, para os Componentes de Formação Geral ( seções 6.2.1 e 6.2.2 ) e de Conhecimento Específico ( seções 6.2.3 e 6.2.4 ).
De cada componente, uma das questões foi escolhida para exemplificar a análise gráfica, relacionando as alternativas escolhidas pelos estudantes ( inclusive o gabarito ) ao número de acertos no componente.
O Anexo I apresenta a íntegra da análise gráfica para todas as questões objetivas.
Para cada uma das questões discursivas, os conteúdos dos tipos mais comuns de respostas dos estudantes são apresentados e comparados com o padrão de respostas esperado ( ver Anexo VIII com o padrão de respostas ).
Tomando -se como base as duas questões discursivas do Componente de Formação Geral, a seção 6.3.1.6 apresenta comentários sobre a correção das respostas em relação à Língua Portuguesa.
As tabelas apresentam as seguintes estatísticas das notas20 : média do desempenho na prova, erro padrão da média, desvio padrão, nota mínima, mediana e nota máxima para cada um de seus componentes.
Tais estatísticas contemplam o total de estudantes concluintes da Área de Ciências Biológicas ( Licenciatura/Bacharelado ) inscritos e presentes à prova do Enade/2017, tendo em vista agregações, ou por Grandes Regiões, e o país como um todo, ou por Categoria Administrativa, Organização Acadêmica da IES e Modalidade de Ensino.
Em relação aos gráficos de distribuição de notas, o intervalo considerado foi de 10 unidades, aberto à esquerda e fechado à direita, com exceção do primeiro intervalo, [ 0 ; 10 ], fechado em ambos os extremos.
Para os gráficos de distribuição das notas das questões discursivas, foram consideradas mais duas categorias : questão em branco21 e nota zero. 20 Uma definição dessas estatísticas pode ser encontrada no Glossário. 21 Nesse grupo, estão incluídas também as respostas classificadas como nulas ou desconsideradas. 6.1 ESTATÍSTICAS BÁSICAS DA PROVA Esta seção apresenta estatísticas selecionadas e histogramas da nota geral ( 6.1.1 ) e de cada componente : Formação Geral ( 6.1.2 ) e Conhecimento Específico ( 6.1.3 ).
São também apresentadas estatísticas selecionadas de subpopulações caracterizadas por Grande Região, Categoria Administrativa, Organização Acadêmica e Modalidade de Ensino. 6.1.1 Estatísticas Básicas Gerais A Tabela 6.1 apresenta as Estatísticas Básicas da prova, por Grande Região segundo Modalidade de Ensino e Habilitação, dos estudantes concluintes de Ciências Biológicas ( Licenciatura/Bacharelado ).
As análises serão feitas primeiramente para os cursos de Licenciatura e depois para os de Bacharelado.
Cada Habilitação é desagregada internamente em cursos Presenciais e cursos a Distância.
A Média das notas da prova como um todo dos alunos de Licenciatura ( nas seções seguintes serão analisados os Componentes de Formação Geral e de Conhecimento Específico ) foi 44,4, sendo que os alunos da região Norte obtiveram a Média mais baixa ( 40,7 ), e os da região Sudeste, a Média mais alta ( 47,7 ).
As demais médias foram : 42,3 na região Nordeste, 44,7 na região Sul e 43,1 na região Centro-Oeste.
Já a Média das notas para os alunos de cursos Presenciais de Licenciatura foi 45,1, sendo que os alunos da região Norte obtiveram a Média mais baixa ( 40,7 ), e os da região Sul obtiveram a Média mais alta ( 48,5 ).
As demais médias foram : 42,8 na região Nordeste, 48,3 na região Sudeste e 43,7 na região Centro-Oeste.
O Desvio padrão para o Brasil como um todo foi 13,5, sendo o maior Desvio padrão encontrado na região Sudeste ( 13,8 ), e o menor, nas regiões Norte e Nordeste ( 12,8 ), indicando uma dispersão um pouco menor das notas dessas últimas regiões.
A região que obteve a maior nota Máxima foi a Centro-Oeste ( 90,3 ), ao passo que a região que atingiu a menor nota Máxima foi a Norte ( 80,1 ).
A Mediana do Brasil, como um todo, foi 45,0, sendo a maior Mediana obtida na região Sul ( 49,2 ), e a menor obtida na Norte ( 40,6 ).
A nota Mínima foi zero em todas as regiões, com exceção da região Centro-Oeste, com 10,6.
Considerando -se as notas segundo Grande Região, observa -se que existe diferença estatisticamente significativa ao nível de 95 \% entre as maiores Médias, obtidas nas regiões Sul ( 48,5 ) e Sudeste ( 48,3 ), e as médias de todas as demais regiões²². ²² Todas as comparações deste capítulo utilizam os intervalos de 95 \%.
Os erros-padrões da média que possibilitam os testes estão disponíveis nas tabelas.
Já a Média das notas da prova dos alunos de cursos a Distância de Licenciatura, 41,2, foi mais baixa do que a dos alunos de cursos Presenciais, sendo que os alunos da região Centro-Oeste obtiveram a Média mais baixa ( 35,3 ), e os da região Sudeste obtiveram a Média mais alta ( 45,6 ).
As demais médias foram : 37,8 nas regiões Norte e Sul, e 37,5 na região Nordeste.
O Desvio padrão para o Brasil, como um todo, foi 13,8, sendo o maior Desvio padrão encontrado na região Centro-Sudeste ( 14,0 ), e o menor, na região Centro-Oeste ( 10,6 ).
A região que obteve a maior nota Máxima foi a Sudeste ( 88,4 ), ao passo que a região que atingiu a menor nota Máxima foi a Centro-Oeste ( 66,7 ).
A Mediana do Brasil, como um todo, foi 40,9, sendo a maior Mediana obtida na região Sudeste ( 46,0 ), e a menor, obtida na região Norte ( 34,4 ).
A nota Mínima foi zero nas regiões Sudeste e Sul.
Considerando -se as notas, segundo Grande Região, observa -se que existe diferença estatisticamente significativa ao nível de 95 \% entre a maior Média, obtida na região Sudeste ( 45,6 ), e as médias das demais regiões.
Cumpre notar que ainda que as notas dos cursos Presenciais sejam comparáveis às dos cursos a Distância dentro de uma mesma Habilitação, não é possível comparar as notas obtidas pelos alunos de Licenciatura e Bacharelado, já que as provas são, em parte, diferentes.
Não foram avaliados cursos a Distância de Bacharelado.
A Média das notas da prova, como um todo, para os alunos de cursos Presenciais de Bacharelado, foi 50,7, sendo que os alunos da região Norte obtiveram a Média mais baixa ( 44,3 ), e os da região Sudeste obtiveram a Média mais alta ( 51,7 ).
As demais médias foram : 49,6 na região Nordeste, 51,0 na região Sul e 49,6 na região Centro-Oeste.
A região que obteve a maior nota Máxima foi a Sudeste ( 91,5 ) e a com menor nota Máxima foi a Norte ( 81,6 ).
A Mediana do Brasil, como um todo, foi 51,0, sendo a maior Mediana obtida na região Sudeste ( 52,1 ), e a menor obtida na Norte ( 44,8 ).
A nota Mínima foi zero em duas das regiões.
As exceções foram a região Norte, com 14,1, a região Sul, com 10,6 e a região Centro-Oeste, com 12,9.
Considerando -se as notas, segundo Grande Região, observa -se que existe diferença estatisticamente significativa ao nível de 95 \% entre a menor Média, obtida na região Norte ( 44,3 ), e as médias de todas as demais regiões.
O comportamento das notas dos estudantes de todo o Brasil pode ser observado no Gráfico 6.1, que apresenta um histograma com a distribuição das mesmas, segundo a Habilitação.
Essa é uma distribuição unimodal para ambas as Habilitações, com as modas nos intervalos ( 40 ; 50 ] e ( 50 ; 60 ], respectivamente para os alunos de cursos de Licenciatura e de Bacharelado.
A Tabela 6.2 apresenta informações referentes às Estatísticas Básicas da prova desagregadas por Categoria Administrativa e por Organização Acadêmica, segundo a Habilitação.
Os estudantes de Licenciatura de IES Públicas obtiveram Média mais alta ( 45,3 ), e os das IES Privadas obtiveram Média mais baixa ( 42,8 ) que a Média nacional para essa habilitação ( 44,4 ).
Observa -se que existe diferença estatisticamente significativa entre as médias das notas das IES Públicas e Privadas.
O Desvio padrão para as IES Privadas ( 13,5 ) foi inferior ao do Brasil, como um todo ( 13,7 ), indicando uma dispersão um pouco menor das notas nessa Categoria Administrativa do que para o conjunto de alunos de cursos de Licenciatura.
No tocante à Organização Acadêmica, alunos de Licenciatura de Universidades obtiveram nota Média mais alta ( 45,8 ) do que os de Centros Universitários ( 42,7 ), que por sua vez obtiveram nota média mais alta do que os de CEFET/IFET ( 42,3 ) e os de Faculdades ( 39,1 ), diferenças estatisticamente significativas ao nível de 95 \% entre a Média mais alta e as demais médias, assim como entre a Média mais baixa e as demais.
Já em relação à Categoria Administrativa, os estudantes de Bacharelado das IES Públicas obtiveram Média mais alta ( 53,1 ), e os das IES Privadas obtiveram Média mais baixa ( 47,8 ) que a Média nacional para essa habilitação ( 50,7 ).
Observa -se que existe diferença estatisticamente significativa entre as médias das notas das IES Públicas e Privadas.
O Desvio padrão para as IES Privadas ( 12,6 ) foi inferior ao do Brasil, como um todo ( 13,2 ), indicando uma dispersão menor das notas nessa Categoria Administrativa do que para o conjunto de alunos de cursos de Bacharelado.
No tocante à Organização Acadêmica, alunos de CEFET/IFET obtiveram nota mais alta ( 53,5 ) do que os de Universidades ( 51,8 ), que por sua vez obtiveram nota média mais alta do que os de Centros Universitários ( 47,0 ) e os de Faculdades ( 44,5 ), diferença estatisticamente significativa ao nível de 95 \% entre a média mais baixa e as demais médias. 6.1.2 Estatísticas Básicas no Componente de Formação Geral A Tabela 6.3 apresenta as Estatísticas Básicas em relação ao componente da prova que avalia a Formação Geral dos estudantes concluintes, por Grande Região segundo Habilitação e Modalidade de Ensino.
Os alunos de cursos de Licenciatura obtiveram desempenho médio de 51,9.
Quanto à variabilidade, o Desvio padrão das notas dos estudantes do Brasil, como um todo, foi 16,3.
A maior Média foi obtida na região Sudeste ( 54,2 ), e a menor, na região Norte ( 48,7 ).
As demais médias foram : 50,2 na região Nordeste, 52,7 na região Sul e 51,3 na região Centro-Oeste.
Já os alunos de cursos Presenciais de Licenciatura obtiveram desempenho médio de 52,6.
Quanto à variabilidade, o Desvio padrão das notas dos estudantes do Brasil, como um todo, foi 16,2.
A maior Média foi obtida na região Sul ( 56,0 ), e a menor, na região Norte ( 48,8 ).
As demais médias foram : 50,6 na região Nordeste, 54,8 na região Sudeste e 51,9 na região Centro-Oeste.
Já o maior Desvio padrão foi obtido na região Norte ( 16,4 ), e o menor, na região Sul ( 15,5 ).
Os demais desvios padrões foram : 15,8 na região Nordeste, 16,3 na região Sudeste e 16,1 na região Centro-Oeste.
A maior nota no Componente de Formação Geral da prova de alunos de cursos Presenciais de Licenciatura foi 98,4, obtida por, pelo menos, um aluno na região Nordeste.
A menor nota Máxima foi obtida na região Centro-Oeste ( 93,0 ).
A Mediana do Brasil, como um todo, foi 53,6, sendo a menor Mediana encontrada na região Norte ( 50,0 ), e a maior, na região Sul ( 57,7 ).
A nota Mínima nessa parte foi zero, em todas as regiões, com exceção da região Centro-Oeste, com 7,5.
Considerando -se as notas segundo Grande Região, observa -se que existe diferença estatisticamente significativa entre a menor Média das notas do Componente de Formação Geral, obtidas na região Norte ( 48,8 ) e as demais médias.
Já a Média das notas do Componente de Formação Geral dos alunos de cursos de Licenciatura a Distância, foi 48,8, mais baixa do que a dos alunos de cursos Presenciais, sendo que os alunos da região Centro-Oeste obtiveram a Média mais baixa ( 43,3 ), e os da região Sudeste obtiveram a Média mais alta ( 51,9 ).
As demais médias foram : 46,0 na região Norte, 45,9 na região Nordeste e 46,7 na região Sul.
O Desvio padrão para o Brasil, como um todo, foi 16,7, sendo o maior Desvio padrão encontrado na região Sudeste ( 17,1 ) e o menor, na região Centro-Oeste ( 12,1 ).
A região que obteve a maior nota Máxima foi a Sudeste ( 95,0 ), ao passo que a região que atingiu a menor nota Máxima foi a Centro-Oeste ( 79,2 ).
A Mediana do Brasil, como um todo, foi 48,9, sendo a maior Mediana obtida na região Sudeste ( 53,0 ), e a menor obtida na região Norte ( 41,0 ).
A nota Mínima foi 21,6 na região Norte, 7,6 na região Centro-Oeste e zero nas demais regiões.
Considerando -se as notas, segundo Grande Região, observa -se que existe diferença estatisticamente significativa ao nível de 95 \% entre a menor Média, obtida nas região Centro-Oeste ( 43,3 ), e a maior Média, obtida na região Sudeste ( 51,9 ).
Como já comentado, não foram avaliados cursos a Distância de Bacharelado.
A Média das notas dos alunos de cursos Presenciais de Bacharelado, foi 56,9, sendo que os alunos da região Norte obtiveram a Média mais baixa ( 52,9 ), e os da região Sudeste obtiveram a Média mais alta ( 57,5 ).
As demais médias foram : 56,2 na região Nordeste, 57,1 na região Sul e 56,6 na região Centro-Oeste.
A região que obteve a maior nota Máxima foi a Sudeste ( 97,8 ) e a com menor nota Máxima foi a Norte ( 89,3 ).
A Mediana do Brasil, como um todo, foi 58,0, sendo a maior Mediana obtida na região Sul ( 58,5 ), e a menor obtida na região Norte ( 52,9 ).
A nota Mínima foi zero em duas das regiões.
As exceções foram a região Norte, com 13,2, e as regiões Sul e Centro-Oeste, ambas com 7,5.
Considerando -se as notas segundo Grande Região, observa -se que existem diferenças estatisticamente significativas ao nível de 95 \% entre a menor Média, obtida na região Norte ( 52,9 ), e as demais médias.
Norte ( 52,9 ).
A nota Mínima foi zero em duas das regiões.
As exceções foram a região Norte, com 13,2, e as regiões Sul e Centro-Oeste, ambas com 7,5.
Considerando -se as notas segundo Grande Região, observa -se que existem diferenças estatisticamente significativas ao nível de 95 \% entre a menor Média, obtida na região Norte ( 52,9 ), e as demais médias.
O Gráfico 6.2 propicia a avaliação do desempenho dos estudantes do Componente de Formação Geral, a partir do histograma da distribuição das notas correspondentes.
A distribuição é unimodal, com moda em ( 50 ; 60 ] para ambas as Habilitações.
Esse foi também o intervalo modal da distribuição de notas da prova, como um todo, para os alunos de Bacharelado ( Gráfico 6.1 ).
Nota -se, ainda, que no Gráfico 6.2, as notas apresentam dispersão maior, quando comparadas às do Gráfico 6.1 ( distribuição das notas da prova ), confirmado pela comparação dos desvios padrões.
Na Tabela 6.4 são apresentadas as informações referentes ao desempenho dos concluintes do Componente de Formação Geral, em diferentes agregações : Categoria Administrativa, Organização Acadêmica e Habilitação.
Primeiramente, com respeito aos alunos de Licenciatura, observa -se que existe diferença estatisticamente significativa entre as médias dos tipos de Categoria Administrativa : a maior Média obtida por estudantes de IES Públicas ( 52,7 ) e a menor pelos de IES Privadas ( 50,3 ).
A ordenação por Organização Acadêmica foi : Universidades ( 53,0 ), CEFET/IFET e Centros Universitários ( 50,5 em ambas ) e Faculdades ( 47,1 ).
A diferença entre a Média das Universidades e a Média das demais organizações acadêmicas é estatisticamente significativa, ao nível de 95 \%.
Dentre os alunos de Bacharelado, observa -se que existe diferença estatisticamente significativa entre as médias dos tipos de Categoria Administrativa : a maior Média obtida por estudantes de IES Públicas ( 59,2 ), e a menor, pelos de IES Privadas ( 54,2 ).
A ordenação por Organização Acadêmica foi : CEFET/IFET ( 58,9 ), Universidades ( 57,8 ), Centros Universitários ( 53,7 ) e Faculdades ( 52,0 ).
A diferença entre as duas maiores médias ( CEFET/IFET e Universidades ) e as duas menores ( Centros universitários e Faculdades ) é estatisticamente significativa, ao nível de 95 \%. 6.1.3 Estatísticas Básicas do Componente de Conhecimento Específico A Tabela 6.5 apresenta as Estatísticas Básicas referentes ao Componente de Conhecimento Específico da área de Ciências Biológicas ( Licenciatura/Bacharelado ), por Grande Região, segundo Habilitação e Modalidade de Ensino.
Quanto ao desempenho de Conhecimento Específico dos alunos dos cursos de Licenciatura, a Média foi 41,9.
A maior Média foi obtida na região Sudeste ( 45,5 ), e a menor, na região Norte ( 38,0 ).
As demais médias foram : 39,6 na região Nordeste, 42,0 na região Sul e 40,3 na região Centro-Oeste.
Já para os alunos dos cursos Presenciais, a Média do desempenho de Conhecimento Específico foi maior, 42,6.
A maior Média foi obtida na região Sudeste ( 46,1 ), e a menor, na região Norte ( 38,0 ).
As demais médias foram : 40,1 na região Nordeste, 46,0 na região Sul e 41,0 na região Centro-Oeste.
Quanto à variabilidade das notas, o Desvio padrão do Brasil, como um todo, foi 14,8, sendo o maior Desvio padrão observado na região Sudeste ( 15,1 ), e o menor, na região Norte ( 13,8 ).
Os demais desvios foram : 14,1 na região Nordeste, 14,5 na região Sul e 14,2 na região Centro-Oeste.
A Mediana das notas dos estudantes de cursos Presenciais de Licenciatura de todo o Brasil foi 42,5.
A maior Mediana ocorreu nas regiões Sudeste e Sul ( 46,3 ), e a menor, na região Norte ( 38,7 ).
As demais medianas foram : 39,7 na região Nordeste e 39,8 na região Centro-Oeste.
A nota Máxima do Brasil, como um todo, foi 96,5, sendo obtida por, pelo menos, um aluno na região Sudeste.
As demais notas máximas foram : 76,8 na região Norte, 83,5 na região Nordeste, 88,5 na região Sul e 94,2 na região Centro-Oeste.
A nota Mínima foi zero em todas as regiões, com exceção da região Centro-Oeste com 7,7.
Observa -se que existe diferença estatisticamente significativa das duas maiores Médias das notas do Componente de Conhecimento Específico ( Sudeste e Sul ) em relação às demais regiões ( Norte, Nordeste e Centro-Oeste ).
Já a Média do desempenho de Conhecimento Específico dos alunos dos cursos a Distância de Licenciatura foi menor, 38,6.
A maior Média foi obtida na região Sudeste ( 43,5 ), e a menor, na região Centro-Oeste ( 32,6 ).
As demais médias foram : 35,1 na região Norte, 34,7 na região Nordeste e 34,8 na região Sul.
O Desvio padrão de todos os alunos de cursos a Distância de Licenciatura foi 15,0, sendo o maior Desvio padrão observado na região Sudeste ( 15,0 ), e o menor, na região Norte ( 11,4 ).
Os demais desvios foram : 14,4 na região Nordeste, 13,6 na região Sul e 12,2 na região Centro-Oeste.
A Mediana das notas dos estudantes de cursos a Distância de Licenciatura foi 38,7.
A maior Mediana ocorreu na região Sudeste ( 43,8 ), e a menor, na região Centro-Oeste ( 31,0 ).
As demais medianas foram : 34,0 nas regiões Norte e Nordeste e 33,9 na região Sul.
A nota Máxima do Brasil, como um todo, foi 89,8, sendo obtida por, pelo menos, um aluno na região Sudeste.
As demais notas máximas foram : 62,6 na região Norte, 86,3 na região Nordeste, 81,0 na região Sul e 73,4 na região Centro-Oeste.
A nota Mínima foi zero nas regiões Sudeste e Sul.
Observa -se que existe diferença estatisticamente significativa da maior Média das notas do Componente de Conhecimento Específico, obtida na região Sudeste ( 43,5 ), em relação às demais regiões.
Para os alunos de cursos Presenciais de Bacharelado, a Média do desempenho de Conhecimento foi 48,6.
A maior Média foi obtida na região Sudeste ( 49,8 ), e a menor, na região Norte ( 41,5 ).
As demais médias foram : 47,3 na região Nordeste, 48,9 na região Sul e 47,3 na região Centro-Oeste.
O Desvio padrão de todos os alunos de cursos Presenciais de Bacharelado foi 14,7, sendo o maior Desvio padrão observado nas regiões Nordeste e Centro-Oeste ( 14,8 ), e o menor, na região Norte ( 14,2 ).
Os demais desvios foram : 14,6 na região Sudeste e 14,4 na região Sul.
A Mediana das notas dos estudantes de cursos Presenciais de Bacharelado foi 48,8.
A maior Mediana ocorreu na região Sudeste ( 50,0 ), e a menor, na região Norte ( 42,1 ).
As demais medianas foram : 47,8 na região Nordeste, 48,8 na região Sul e 47,5 na região Centro-Oeste.
A nota Máxima do Brasil, como um todo, foi 94,0, obtida por pelo menos, um aluno na região Sudeste.
As demais notas máximas foram : 81,3 na região Norte, 92,2 na região Nordeste, 90,2 na região Sul e 92,4 na região Centro-Oeste.
A nota Mínima foi 0,0 em duas regiões, Nordeste e Sudeste.
Na região Norte a nota Mínima foi 11,6, na região Sul foi 4,5 e na região Centro-Oeste, 4,1.
Observa -se que para os alunos de cursos Presenciais de Bacharelado, existe uma diferença estatisticamente significativa da menor Média das notas do Componente de Conhecimento Específico, obtida na região Norte, em relação às demais regiões.
Assim como os Gráficos 6.1 e 6.2, o Gráfico 6.3, apresentado a seguir, permite uma avaliação do desempenho de concluintes de Ciências Biológicas ( Licenciatura/Bacharelado ) segundo a Habilitação, em relação ao Componente de Conhecimento Específico com um histograma da distribuição das notas correspondentes.
Esta também é uma distribuição unimodal para ambas as habilitações, e o grupo modal é, respectivamente, ( 30 ; 40 ] e ( 40 ; 50 ], para estudantes de Licenciatura e Bacharelado.
A Tabela 6.6 apresenta uma comparação dos resultados em relação à Categoria Administrativa e à Organização Acadêmica, segundo a Habilitação, agora levando -se em conta o desempenho de alunos do Componente de Conhecimento Específico da prova.
No que se refere à Organização Acadêmica, para os alunos de Licenciatura, a maior Média foi a de Universidades ( 43,3 ), vindo a seguir a dos Centros Universitários ( 40,1 ), a de CEFET/IFET ( 39,6 ) e por último, a das Faculdades ( 36,4 ).
O maior Desvio padrão, e acima do valor para o Brasil como um todo ( 14,9 ), foi a dos Centros Universitários ( 15,0 ).
As Universidades obtiveram a maior nota Máxima, igual a 96,5, com os Centros Universitários em segundo lugar, com 89,8.
As Faculdades obtiveram nota Máxima igual a 86,3 e os CEFET/IFET, 85,2.
As medianas foram : 42,8 nas Universidades, a maior, 39,4 nos Centros Universitários, 35,4 nas Faculdades, a menor delas, e 38,8 nos CEFET/IFET.
A nota Mínima foi zero em todas as Organizações acadêmicas, com exceção dos CEFET/IFET ( 5,8 ).
Observa -se que existem diferenças estatisticamente significativas ao nível de 95 \% no Componente de Conhecimento Específico entre a maior nota Média, a das Universidades, em relação às Médias das demais Organizações Acadêmicas.
Quanto à Categoria Administrativa, observa -se um comportamento semelhante àquele da parte de Formação Geral e à prova como um todo, ou seja, existe diferença estatisticamente significativa entre as médias das IES Públicas ( 42,7 ) e IES Privadas ( 40,3 ).
Nesse caso também, a maior Média foi obtida por alunos de IES Públicas de ensino.
Para os alunos de Bacharelado, a maior Média, considerando -se a Organização Acadêmica, foi a de CEFET/IFET ( 51,6 ), e a menor, a das Faculdades ( 41,9 ).
As Universidades obtiveram a maior nota Máxima, igual a 94,0.
As medianas foram : 50,0 nas Universidades, 44,5 nos Centros Universitários, 41,7 nas Faculdades, a menor delas, e 53,8 nos CEFET/IFET, a maior delas.
A nota Mínima foi zero para Universidades e para Centros Universitários.
Observa -se que existem diferenças estatisticamente significativas ao nível de 95 \% no Componente de Conhecimento Específico entre a menor nota Média, a das Faculdades, em relação às Médias das demais Organizações Acadêmicas.
Quanto à Categoria Administrativa, observa -se um comportamento semelhante àquele da parte de Formação Geral e à prova, como um todo, para os alunos de Bacharelado, ou seja, existe diferença estatisticamente significativa entre as médias das IES Públicas ( 51,1 ) e as das IES Privadas ( 45,7 ).
Nesse caso, também, a maior Média foi obtida por alunos de IES Públicas de ensino. 6.2 ANÁLISE DAS QUESTÕES OBJETIVAS Esta seção apresenta estatísticas selecionadas e histogramas das Questões Objetivas de Formação Geral para Licenciatura ( 6.2.1 ) e para Bacharelado ( 6.2.2 ), e das Questões Objetivas de Conhecimento Específico para Licenciatura ( 6.2.3 ) e para Bacharelado ( 6.2.4 ).
São também apresentadas e comparadas as médias das subpopulações caracterizadas por Grande Região. 6.2.1 Componente de Formação Geral – Ciências Biológicas ( Licenciatura ) A Tabela 6.7a apresenta as Estatísticas Básicas relativas às oito questões objetivas do componente da prova que abrange a Formação Geral dos estudantes de Ciências Biológicas ( Licenciatura ), segundo a Modalidade de Ensino.
A Média dos alunos nestas questões foi 48,3.
A menor Média foi encontrada na região Norte ( 44,3 ), e a maior, na região Sudeste ( 50,8 ).
As demais médias foram : 46,4 na região Nordeste, 49,3 na região Sul e 48,1 na região Centro-Oeste.
Para os alunos de cursos Presenciais, a Média foi 49,0.
A menor Média foi encontrada na região Norte ( 44,4 ), e a maior, na região Sul ( 53,0 ).
As demais médias foram : 46,7 na região Nordeste, 51,4 na região Sudeste e 48,9 na região Centro-Oeste.
O Desvio padrão do Brasil foi 19,5, sendo o maior Desvio padrão encontrado nas regiões Sudeste e Centro-Oeste ( 19,7 em ambos ), e o menor, na região Nordeste ( 18,8 ).
Os demais desvios foram : 19,4 na região Norte e 19,1 na região Sul.
As medianas do Brasil como um todo e todas as cinco regiões foi 50,0.
A nota Máxima 100,0 foi alcançada em todas as grandes Regiões.
As notas mínimas ( 0,0 ) foram iguais para todas as regiões.
Já para os alunos de cursos a Distância, a Média foi menor, 45,1.
A menor Média foi encontrada na região Centro-Oeste ( 38,7 ), e a maior, na região Sudeste ( 48,5 ).
As demais médias foram : 42,2 na região Norte, 42,7 na região Nordeste e 42,6 na região Sul.
O Desvio padrão do Brasil para esta Modalidade de Ensino nos cursos de Licenciatura foi 20,1, sendo o maior Desvio padrão encontrado na região Sudeste ( 21,0 ), e o menor, na região Centro-Oeste ( 15,9 ).
Os demais desvios foram : 19,5 na região Norte, 18,0 na região Nordeste e 19,6 na região Sul.
As medianas do Brasil, como um todo, e de para a região Sudeste foram 50,0.
Para as demais regiões, a mediana foi mais baixa, 37,5.
A nota Máxima ( 100,0 ) foi alcançada em quase todas as grandes Regiões, com exceção da região Norte ( 87,5 ) e da Região Centro-Oeste ( 75,0 ).
As notas mínimas ( 0,0 ) foram as mesmas para todas as regiões.
A Tabela 6.8a apresenta o Índice de Facilidade e o Índice de Discriminação ( ponto-bisserial ) para cada uma das questões objetivas do Componente de Formação Geral.
Quanto ao Índice de Facilidade, foram usadas as seguintes cores para diferenciar o nível de dificuldade da questão :  Azul para as questões classificadas com índice Muito fácil ( \&gt ; =0,86 ), verde para as questões classificadas com índice Fácil ( 0,61 a 0,85 ), amarelo para as questões classificadas com Médio ( 0,41 a 0,60 ), vermelho para as questões classificadas com Difícil ( 0,16 a 0,40 ) e roxo para as questões classificadas com Muito difícil ( \&lt ; =0,15 ).
Já quanto ao Índice de Discriminação, foram usadas as seguintes cores para qualificar a questão :  As questões classificadas com índice Fraco receberam a cor vermelha ( \&lt ; =0,19 ), as classificadas com Médio receberam a cor amarela ( 0,20 a 0,29 ), as classificadas com Bom receberam a cor verde ( 0,30 a 0,39 ), e as classificadas com Muito bom ( \&gt ; =0,40 ) receberam a cor azul.
As questões objetivas do Componente de Formação Geral, segundo o Índice de Facilidade, foram assim avaliadas : das oito questões, nenhuma teve o Índice de Facilidade classificado como Muito fácil, e duas questões foram tidas como Fácil.
Três questões foram consideradas com índice de dificuldade Médio, situando -se no intervalo entre 0,41 e 0,60 do A Tabela 6.8a apresenta o Índice de Facilidade e o Índice de Discriminação ( ponto-bisserial ) para cada uma das questões objetivas do Componente de Formação Geral.
Quanto ao Índice de Facilidade, foram usadas as seguintes cores para diferenciar o nível de dificuldade da questão :  Azul para as questões classificadas com índice Muito fácil ( \&gt ; =0,86 ), verde para as questões classificadas com índice Fácil ( 0,61 a 0,85 ), amarelo para as questões classificadas com Médio ( 0,41 a 0,60 ), vermelho para as questões classificadas com Difícil ( 0,16 a 0,40 ) e roxo para as questões classificadas com Muito difícil ( \&lt ; =0,15 ).
Já quanto ao Índice de Discriminação, foram usadas as seguintes cores para qualificar a questão :  As questões classificadas com índice Fraco receberam a cor vermelha ( \&lt ; =0,19 ), as classificadas com Médio receberam a cor amarela ( 0,20 a 0,29 ), as classificadas com Bom receberam a cor verde ( 0,30 a 0,39 ), e as classificadas com Muito bom ( \&gt ; =0,40 ) receberam a cor azul.
As questões objetivas do Componente de Formação Geral, segundo o Índice de Facilidade, foram assim avaliadas : das oito questões, nenhuma teve o Índice de Facilidade classificado como Muito fácil, e duas questões foram tidas como Fácil.
Três questões foram consideradas com índice de dificuldade Médio, situando -se no intervalo entre 0,41 e 0,60 do Índice de Facilidade, ou seja, houve entre 41,0 \% e 60,0 \% de acertos, enquanto duas questões foram classificadas na categoria Difícil, situando -se no intervalo entre 0,16 e 0,40.
Por fim, uma das questões apresentou menos de 15 \% de acertos, razão pela qual foi classificada como Muito difícil.
O Índice de Facilidade variou de 0,12 a 0,84, e o de Discriminação, de 0,30 a 0,52.
A Tabela 6.9a apresenta a distribuição das questões considerando -se simultaneamente a classificação dos dois índices.
As seis questões com Índice de Discriminação Muito bom figuraram entre três níveis de dificuldade : Fácil, Médio e Difícil.
Dessas, duas foram classificadas na categoria Fácil ( questões 4 e 6 ) do Índice de Facilidade, três, na categoria Médio ( questões 2, 5 e 8 ), e uma foi classificada na categoria Difícil ( questão 7 ).
Em particular, a questão 4 foi a que apresentou o maior poder discriminatório, com índice 0,52, tendo sido considerada Fácil em termos de facilidade, com uma proporção de 0,66 acertos.
O máximo de acertos foi alcançado pela questão 6, com um Índice de Facilidade de 0,84.
A questão de número 1 apresentou um Índice de Facilidade de 0,12 ( Muito difícil ), ou seja, um pouco mais de um décimo dos estudantes conseguiu resolvê -la, dentro do universo de participantes.
Seu Índice de Discriminação foi Bom ( 0,30 ).
O Gráfico 6.4a, para exemplificar, analisa o comportamento da questão de número 4 de Formação Geral.
Trata -se de uma questão considerada Fácil, com relação à facilidade, e foi a que obteve o segundo maior Índice de Discriminação dessa parte da prova.
Neste gráfico, cada uma das cinco curvas representa o percentual de respostas em determinada alternativa da questão, em função do número de acertos dos estudantes nessa parte da prova ( Formação Geral/Múltipla Escolha ), antes de possíveis eliminações pelo critério do ponto-bisserial.
A curva em verde corresponde à alternativa B, a alternativa correta para essa questão.
Observa -se que entre os estudantes com menor número de acertos, nessa parte do exame, a situação mais frequente foi a escolha da alternativa E ( em vermelho ), incorreta.
Por exemplo, entre os estudantes que acertaram três questões, 21,8 \% escolheram a alternativa E, 14,8 \% escolheram a alternativa D ( roxo ), mais da metade ( 53,7 \% ) escolheu a alternativa B ( em verde ), 6,5 \% escolheram a alternativa A ( azul ) e 3,2 \% escolheram a alternativa C ( preto ).
Entre os que acertaram três respostas entre as questões de múltipla escolha de Formação Geral, nenhum deixou a questão em branco ou marcou mais de uma alternativa, invalidando a questão.
À medida que o número de acertos aumenta, indicando desempenho melhor nessa parte da prova, aumenta concomitantemente a proporção de estudantes que selecionaram a alternativa correta B, atingindo 100 \% para os estudantes com oito acertos.
Essa análise permite verificar como a questão discriminou os grupos de desempenho, justificando o alto índice obtido na questão.
Cumpre notar que não é possível inferir deste gráfico nem o índice de facilidade ( que seria uma média da proporção ponderada pela quantidade de alunos com cada uma das notas ), nem o índice de discriminação ponto-bisserial, por razão equivalente.
No caso extremo, no qual a grande concentração dos acertos dos alunos fosse abaixo de quatro, o índice de facilidade seria obrigatoriamente abaixo de 40 \% ( neste exemplo ).
Caso a concentração fosse em seis acertos ou mais, o índice seria obrigatoriamente acima de 60 \%.
Os gráficos relativos às demais questões de Formação Geral constam do Anexo Ia. 6.2.2 Componente de Formação Geral – Ciências Biológicas ( Bacharelado ) A Tabela 6.7b apresenta as Estatísticas Básicas relativas às oito questões objetivas do componente da prova que abrange a Formação Geral dos estudantes de Ciências Biológicas ( Bacharelado ), segundo a Modalidade de Ensino.
Como já comentado, não foram observados cursos a Distância de Bacharelado.
Para os alunos de cursos Presenciais de Bacharelado, a Média foi 54,5.
A menor Média foi encontrada na região Norte ( 48,5 ), e a maior, na região Sul ( 56,0 ).
As demais médias foram : 53,6 na região Nordeste, 54,9 na região Sudeste e 53,9 na região Centro-Oeste.
O Desvio padrão do Brasil foi 18,9, sendo o maior Desvio padrão encontrado na região Sul ( 19,2 ), e o menor, na região Norte ( 18,4 ).
Os demais desvios foram : 18,7 nas regiões Nordeste e Centro-Oeste, e 18,8 na região Sudeste.
As medianas do Brasil como um todo e das cinco regiões foram 50,0.
A nota Máxima ( 100,0 ) foi alcançada em todas as grandes Regiões, com exceção da região Norte ( 87,5 ).
As notas mínimas ( 0,0 ) foram iguais para todas as regiões, com exceção da região Centro-Oeste, que obteve nota Mínima igual a 12,5.
A Tabela 6.8b apresenta o Índice de Facilidade e o Índice de Discriminação ( ponto-bisserial ) para cada uma das questões objetivas do Componente de Formação Geral.
A convenção das cores é a mesma da Tabela 6.8a.
As questões objetivas do Componente de Formação Geral, segundo o Índice de Facilidade, foram assim avaliadas : das oito questões, uma teve o Índice de Facilidade classificado como Muito fácil, e duas questões foram tidas como Fácil.
Duas questões foram consideradas com índice de dificuldade Médio, enquanto três questões foram classificadas na categoria Difícil, situando -se no intervalo entre 0,16 e 0,40.
Nenhuma das questões apresentou menos de 15 \% de acertos, e não foram classificadas como Muito difícil.
O Índice de Facilidade variou de 0,16 a 0,92, e o de Discriminação, de 0,36 a 0,48.
A Tabela 6.9b apresenta a distribuição das questões, considerando -se simultaneamente a classificação dos dois índices, semelhantemente ao apresentado na Tabela 6.9a, mas aqui para os alunos de Bacharelado.
Tabela 6.9a, mas aqui para os alunos de Bacharelado.
As seis questões com Índice de Discriminação Muito bom figuraram entre três níveis de dificuldade : Fácil, Médio e Difícil.
Dessas, duas foram classificadas na categoria Fácil ( as questões 4 e 5 ) do Índice de Facilidade, duas na categoria Médio ( questões 2 e 8 ) e, por fim, duas na categoria Difícil ( questões 3 e 7 ).
Em particular, as questões 5 e 8 foram as que apresentaram o maior poder discriminatório, ambas com índice 0,48, e foram, respectivamente, consideradas Fácil e Médio em termos de facilidade, com uma proporção de 0,66 e de 0,55 acertos.
O máximo de acertos foi alcançado pela questão 6, com um Índice de Facilidade de 0,92.
O Gráfico 6.4b, para exemplificar, analisa o comportamento da questão de número 4 de Formação Geral.
Trata -se de uma questão considerada Fácil, em relação à facilidade, e com Índice de Discriminação Muito bom.
Semelhantemente ao apresentado no Gráfico 6.4a, neste gráfico também, cada uma das cinco curvas representa o percentual de respostas em determinada alternativa da questão, em função do número de acertos dos estudantes nessa parte da prova ( Formação Geral/Múltipla Escolha ), antes de possíveis eliminações pelo critério do ponto-bisserial.
A curva em verde corresponde à alternativa B, a alternativa correta para essa questão.
Observa -se que entre os estudantes com menor número de acertos, nessa parte do exame, a situação mais frequente foi a escolha da alternativa E ( em vermelho ), incorreta.
Por exemplo, entre os estudantes que acertaram três questões, 17,2 \% escolheram a alternativa E, mais da metade ( 59,8 \% ) escolheram a alternativa B ( em verde ), 12,9 \% a D ( roxo ), 2,6 \% escolheram a alternativa C ( preto ), e 7,2 \% escolheram a alternativa A ( em azul ).
Entre os que acertaram três respostas entre as questões de múltipla escolha de Formação Geral, 0,2 \% deixou a questão em branco e 0,1 \% marcou mais de uma alternativa, invalidando a questão.
À medida que o número de acertos aumenta, indicando desempenho melhor nessa parte da prova, aumenta concomitantemente a proporção de estudantes que selecionaram a alternativa correta B, atingindo 100 \% para os estudantes com oito acertos.
Essa análise permite verificar como a questão discriminou os grupos de desempenho, justificando o alto índice obtido na questão.
Os gráficos relativos às demais questões de Formação Geral constam do Anexo Ib. 6.2.3 Componente de Conhecimento Específico – Ciências Biológicas ( Licenciatura ) A Tabela 6.10a apresenta as Estatísticas Básicas em relação às questões objetivas do Componente de Conhecimento Específico da prova de Ciências Biológicas ( Licenciatura ) por Grande Região, segundo a Modalidade de Ensino.
Nessa parte da prova, nenhuma das questões foi anulada pela Comissão Assessora da Área.
Como será discutido ainda nesta seção, cinco questões objetivas do Componente de Conhecimentos Específicos não foram usadas no cômputo das notas por terem sido descartadas pelo critério do ponto-bisserial.
Assim, as notas foram calculadas com base em 22 das 27 questões objetivas de Conhecimentos Específicos.
A Média deste componente para os alunos de cursos de Licenciatura foi de 43,1.
A menor Média foi observada na região Norte ( 39,4 ), e a maior, na região Sudeste ( 46,6 ).
O Desvio padrão de todo o Brasil foi 15,8, sendo o menor Desvio padrão encontrado na região Norte ( 14,9 ), e o maior, na região Sul ( 16,3 ).
A situação é um pouco diferente quando há desagregação por Modalidade de Ensino : a Média deste componente para os alunos de cursos Presenciais de Licenciatura foi de 43,9.
A menor Média foi observada na região Norte ( 39,4 ), e a maior, nas regiões Sudeste e Sul ( 47,2 em ambas ).
O Desvio padrão de todo o Brasil foi 15,7, sendo o menor Desvio padrão encontrado na região Norte ( 14,9 ), e o maior, na região Sudeste ( 15,9 ).
A Mediana de todo o Brasil foi 45,5, a mesma das regiões Sudeste e Sul.
A Mediana nas regiões Norte, Nordeste e Centro-Oeste foi menor, 40,9.
A nota Máxima da prova foi 100,0, obtida nas questões objetivas do Componente de Conhecimento Específico, por, pelo menos, um aluno da região Sudeste.
A nota Mínima foi zero em todas as regiões, com exceção da Centro-Oeste ( 9,1 ).
Já para os alunos de cursos a Distância de Licenciatura, a Média deste componente foi mais baixa, 39,7.
A menor Média foi observada na região Centro-Oeste ( 34,3 ), e a maior, na região Sudeste ( 44,7 ).
O Desvio padrão de todo o Brasil foi 16,0, sendo o menor Desvio padrão encontrado na região Norte ( 12,6 ), e o maior, na região Sudeste ( 15,8 ).
A Mediana de todo o Brasil foi 40,9.
Na região Sudeste ( 45,5 ), a Mediana foi maior que a nacional, e nas regiões Norte, Nordeste, Sul ( 36,4 em cada ) e Centro-Oeste ( 31,8 ) foi menor.
A nota Máxima da prova foi 90,9, obtida nas questões objetivas do Componente de Conhecimento Específico, por, pelo menos, um aluno das regiões Nordeste e Sudeste.
A nota Mínima foi zero em quase todas as regiões, com exceção da região Norte ( 13,6 ) e da região Centro-Oeste ( 4,5 ).
A Tabela 6.11a apresenta os Índices de Facilidade e Discriminação ( ponto-bisserial ) das questões objetivas do Componente de Conhecimento Específico para os estudantes de Ciências Biológicas ( Licenciatura ).
Para facilitar a diferenciação das questões, serão usadas as mesmas cores da Tabela 6.8a e 6.8b para as diferentes classificações dos Índices de Facilidade e de Discriminação.
Observa -se que nenhuma das questões foi anulada pela CAA, e que as análises dizem respeito às 27 questões válidas.
A partir do Índice de Facilidade obtido, pode -se concluir que um pouco mais da metade das questões objetivas da prova foi considerada, pelo menos, Difícil : das 27 questões, 14 foram classificadas como Difícil ou como Muito difícil.
Três questões foram classificadas como Fácil, e outras dez consideradas como Médio.
Já quanto aos índices de discriminação das questões objetivas do Componente de Conhecimento Específico da prova, obteve -se como resultado a seguinte classificação : nove das 27 questões foram consideradas como boas, enquanto três delas tiveram Índice de Discriminação Muito bom.
Assim, para um pouco menos da metade das questões – 12 em 27 – os Índices de Discriminação foram Bom ou Muito bom.
Dentre as demais, dez delas foram classificadas como Médio, e outras cinco, como Fraco, sendo 15, por conseguinte, a quantidade de questões nos dois patamares mais baixos de discriminação.
Constata -se, assim, que a prova – no que se refere ao Componente de Conhecimento Específico – possuía capacidade média de discriminar dentre aqueles que dominam ou não o conteúdo.
A Tabela 6.12a apresenta a distribuição das questões, considerando -se simultaneamente a classificação dos dois índices.
Dentre as questões que alcançaram os maiores índices de discriminação, classificadas com Índice de Discriminação Muito bom, uma delas, a de número 21, a com índice mais elevado ( 0,45 ), foi classificada como Médio quanto ao Índice de Facilidade, com 55 \% dos estudantes marcando a opção correta.
A questão de número 18 foi a mais difícil dentre as 27 questões específicas, com baixo Índice de Facilidade, apenas 13 \% de acertos.
Essa questão apresentou poder discriminatório baixo, 0,03, o que comprova ter sido a mais difícil para os estudantes.
Destaca -se, também, a questão 30, com Índice de Facilidade 0,19, o que, em termos percentuais, corresponde a 19 \% de estudantes que responderam acertadamente.
Já 0,19 também foi o seu de Índice de Discriminação.
Tais questões foram, portanto, pelo critério ponto-bisserial, consideradas inadequadas.
Por isso, as questões 18 e 30, além das questões 16, 24 e 25 foram eliminadas do cômputo da nota final.
A título de exemplo das análises do comportamento das questões objetivas, o Gráfico 6.5a analisa a questão 21 do Componente de Conhecimento Específico.
Esta questão foi considerada pelas respostas dos estudantes avaliados na prova como uma questão com nível de facilidade Médio, 0,55, ou seja, 55 \% dos estudantes assinalaram acertadamente a opção E, correspondente ao gabarito.
Como já comentado, seu Índice de Discriminação foi igual a 0,45, classificado como Muito bom, o maior valor de discriminação.
Neste gráfico também, cada uma das cinco curvas representa o percentual de respostas em determinada alternativa da questão 21, em função do número de acertos dos estudantes nessa parte da prova, antes de possíveis eliminações de questões pelo critério do ponto-bisserial.
A alternativa correta E, representada no gráfico pela curva em vermelho, foi escolhida em maiores proporções pelos alunos com desempenho melhor nessa parte da prova.
Já as alternativas incorretas, também denominadas distratores, foram selecionadas, principalmente, por aqueles com notas mais baixas.
Observa -se que a soma não é 100 \%, por causa das questões não respondidas ou com mais de uma opção marcada.
Aqueles com nota zero, na sua quase totalidade, deixaram essa questão em branco ou marcaram mais de uma alternativa, comportamento considerado inválido.
A proporção de alunos que selecionou a resposta correta E aumenta gradativamente, chegando a atingir 100 \% para 21 acertos ou mais, enquanto a proporção dos que escolheram alternativas incorretas decai, a partir de dois ou três acertos, como função do número de acertos nessa parte da prova.
Os gráficos relativos às demais questões do Conhecimento Específico constam do Anexo I. 6.2.4 Componente de Conhecimento Específico – Ciências Biológicas ( Bacharelado ) A Tabela 6.10b apresenta as Estatísticas Básicas em relação às questões objetivas do Componente de Conhecimento Específico da prova de Ciências Biológicas ( Bacharelado ) por Grande Região segundo a Modalidade de Ensino.
Nenhuma das questões foi anulada pela Comissão Assessora da Área.
Porém, como será discutido ainda nesta seção, seis das questões objetivas do Componente de Conhecimentos Específicos não foram usadas no cômputo das notas por terem sido descartadas pelo critério do ponto-bisserial.
Assim, as notas foram calculadas com base em 21 das 27 questões objetivas de Conhecimentos Específicos.
Como já comentado, não foram avaliados cursos a Distância de Bacharelado.
Para os alunos de cursos Presenciais de Bacharelado, a Média do Brasil deste componente foi de 53,1.
A menor Média foi observada na região Norte ( 45,9 ), e a maior, na região Sudeste ( 54,2 ).
O Desvio padrão de todo o Brasil foi 15,9, sendo o menor Desvio padrão encontrado nas regiões Norte e Sul ( 15,7 ), e o maior, na região Nordeste ( 16,1 ).
A Mediana de todo o Brasil foi 52,4, a mesma em todas as regiões, exceto na região Norte ( 47,6 ).
A nota Máxima da prova foi 100,0, obtida nas questões objetivas do Componente de Conhecimento Específico, por, pelo menos, um aluno das regiões Nordeste e Sudeste.
A nota Mínima foi zero em três regiões, com exceção das regiões Norte ( 9,5 ) e Centro-Oeste ( 4,8 ).
A Tabela 6.11b apresenta os Índices de Facilidade e Discriminação ( ponto-bisserial ) das questões objetivas do Componente de Conhecimento Específico, para os estudantes de Ciências Biológicas ( Bacharelado ).
Para facilitar a diferenciação das questões, serão usadas as mesmas convenções definidas para as Tabela 6.8a, 6.8b e 6.11a.
A partir do Índice de Facilidade obtido, pode -se concluir que pouco mais de um terço das questões objetivas da prova foi considerado, pelo menos, Difícil : das 27 questões válidas, dez foram classificadas como Difícil ( nove ) ou como Muito difícil ( uma ).
Sete questões foram classificadas como Fácil, e outras dez foram consideradas como Médio.
Já quanto aos índices de discriminação das questões objetivas do Componente de Conhecimento Específico da prova, obteve -se como resultado a seguinte classificação : 11 das 27 questões foram consideradas como boas, enquanto outras duas tiveram Índice de Discriminação Muito bom.
Assim, para um pouco menos da metade das questões – 13 em 27 – os Índices de Discriminação foram Bom ou Muito bom.
Dentre as demais, oito delas foram classificadas como Médio e seis como Fraco, sendo 14, por conseguinte, a quantidade de questões nos dois patamares mais baixos de discriminação.
Constata -se, assim, que a prova – no que se refere ao Componente de Conhecimento Específico – possuía capacidade média de discriminar dentre aqueles que dominam ou não o conteúdo.
A Tabela 6.12b apresenta a distribuição das questões, considerando -se simultaneamente a classificação dos dois índices.
Dentre as questões que alcançaram os maiores índices de discriminação, classificadas com Índice de Discriminação Muito bom, uma delas, a de número 17, a de índice mais elevado ( 0,44 ), foi classificada como Fácil quanto ao Índice de Facilidade, com 66 \% dos estudantes marcando a opção correta.
A questão de número 18 foi a mais difícil dentre as 27 questões específicas, com baixo Índice de Facilidade, apenas 13 \% de acertos.
Essa questão apresentou poder discriminatório muito baixo, 0,08, o que comprova ter sido a mais difícil para os estudantes.
Tal questão foi, portanto, pelo critério ponto-bisserial, considerada inadequada.
Por isso, a questão 18 foi eliminada do cômputo da nota final, assim como as questões 16, 24, 25, 28 e 33.
A Tabela 6.12b apresenta a distribuição das questões, considerando -se simultaneamente a classificação dos dois índices.
Dentre as questões que alcançaram os maiores índices de discriminação, classificadas com Índice de Discriminação Muito bom, uma delas, a de número 17, a de índice mais elevado ( 0,44 ), foi classificada como Fácil quanto ao Índice de Facilidade, com 66 \% dos estudantes marcando a opção correta.
A questão de número 18 foi a mais difícil dentre as 27 questões específicas, com baixo Índice de Facilidade, apenas 13 \% de acertos.
Essa questão apresentou poder discriminatório muito baixo, 0,08, o que comprova ter sido a mais difícil para os estudantes.
Tal questão foi, portanto, pelo critério ponto-bisserial, considerada inadequada.
Por isso, a questão 18 foi eliminada do cômputo da nota final, assim como as questões 16, 24, 25, 28 e 33.
A título de exemplo das análises do comportamento das questões objetivas, o Gráfico 6.5b analisa a questão 17 do Componente de Conhecimento Específico.
Essa questão foi considerada pelas respostas dos estudantes avaliados na prova como uma questão com nível de facilidade Fácil, 0,66, ou seja, 66 \% dos estudantes assinalaram acertadamente a opção D, correspondente ao gabarito.
Como já comentado, seu Índice de Discriminação foi igual a 0,44, classificado como Muito bom, o maior valor de discriminação.
Neste gráfico, cada uma das cinco curvas representa o percentual de respostas em determinada alternativa da questão 17, em função do número de acertos dos estudantes nessa parte da prova, antes de possíveis eliminações de questões pelo critério do ponto-bisserial.
A alternativa correta D, representada no gráfico pela curva em roxo, foi escolhida em maiores proporções pelos alunos com desempenho melhor nessa parte da prova.
Já as alternativas incorretas, também denominadas distratores, foram selecionadas, principalmente, por aqueles com notas mais baixas.
Observa -se que a soma não é 100 \%, por causa das questões não respondidas ou com mais de uma opção marcada.
Aqueles com nota zero, na sua quase totalidade, deixaram essa questão em branco ou marcaram mais de uma alternativa, comportamento considerado inválido.
A proporção de alunos que selecionou a resposta correta D aumenta gradativamente, chegando a atingir 100 \% para 21, 22 ou 24 acertos, enquanto a proporção dos que escolheram alternativas incorretas decai, a partir de um ou dois acertos, como função do número de acertos nessa parte da prova.
Os gráficos relativos às demais questões do Conhecimento Específico constam do Anexo I. 6.3 ANÁLISE DAS QUESTÕES DISCURSIVAS Esta seção apresenta estatísticas selecionadas e histogramas das Questões Discursivas de Formação Geral ( 6.3.1 ) e Conhecimento Específico ( 6.3.2 ).
São também apresentadas e comparadas as médias de subpopulações, caracterizadas por Grande Região. 6.3.1 Componente de Formação Geral As análises dos resultados de desempenho dos estudantes de Ciências Biológicas ( Licenciatura/Bacharelado ) nas duas questões discursivas relativas à Formação Geral encontram -se na Tabela 6.13 e no Gráfico 6.6, desagregadas por Grande Região, segundo Habilitação e Modalidade de Ensino.
Na Tabela 6.13, observa -se que a nota Média nesse conjunto de questões foi acima da obtida nas objetivas pelos estudantes de Licenciatura.
Esses estudantes obtiveram, em Formação Geral, Média 48,3 nas questões objetivas e 57,4 nas questões discursivas.
A mesma situação acontece para os alunos de Bacharelado : média menor nas objetivas ( 54,5 ) do que nas discursivas ( 60,6 ).
Considerando -se agora os alunos de cursos Presenciais de Licenciatura, observa -se que a nota Média nesse conjunto de questões foi acima da obtida nas objetivas.
Esses estudantes obtiveram, em Formação Geral, Média 49,0 nas questões objetivas e 58,0 nas questões discursivas.
Pode -se notar também que o Desvio padrão nesse conjunto de questões foi um pouco maior do que o obtido nas objetivas, 19,5 nas questões objetivas e 22,7 nas questões discursivas.
A maior Média foi obtida na região Sul ( 60,6 ), e a menor, na região Norte ( 55,4 ).
A Mediana dos alunos de cursos Presenciais de Licenciatura, neste componente, foi 63,0.
Nas regiões Norte ( 60,5 ), Nordeste ( 62,0 ) e Centro-Oeste ( 61,5 ) a Mediana foi menor e nas demais regiões foi maior que o valor para o Brasil : Sudeste ( 64,5 ) e Sul ( 65,0 ).
A nota Máxima ( 98,0 ) foi a mesma na região Sudeste, sendo 96,5 nas regiões Norte e Sul, 97,0 na região Nordeste e 97,5 na região Centro-Oeste.
A nota Mínima ( 0,0 ) foi a mesma em todas as regiões do Brasil.
Também para os alunos de cursos a Distância de Licenciatura, a nota Média nesse conjunto de questões foi acima da obtida nas questões objetivas pelo mesmo grupo de estudantes.
Esses estudantes obtiveram, em Formação Geral, Média 45,1 nas questões objetivas e 54,3 nas questões discursivas.
Pode -se notar que o Desvio padrão nesse conjunto de questões foi também um pouco maior do obtido nas objetivas, 20,1 nas questões objetivas e 22,8 nas questões discursivas.
A maior Média foi obtida na região Sudeste ( 57,1 ), e a menor, na região Centro-Oeste ( 50,3 ).
A Mediana dos alunos de cursos a Distância de Licenciatura, neste componente, foi 59,0.
Em todas as regiões, a Mediana foi menor, exceto na região Sudeste ( 62,5 ).
A nota Máxima ( 99,0 ) foi a mesma na região Sudeste, sendo 92,5 na região Norte, 93,5 na região Nordeste, 97,0 na região Sul e 86,0 na região Centro-Oeste.
A nota Mínima ( 0,0 ) foi a mesma em todas as regiões do Brasil.
Como já comentado, não foram avaliados cursos a Distância de Bacharelado.
Para os alunos de cursos Presenciais de Bacharelado, a nota Média nas questões objetivas também foi menor do que a das discursivas.
Esses estudantes obtiveram nas questões objetivas Média 54,5 e nas questões discursivas, 60,6.
Pode -se notar que o Desvio padrão nesse conjunto de questões foi maior do que o obtido nas objetivas : 18,9 nas questões objetivas e 22,1 nas questões discursivas.
A maior Média foi obtida na região Sudeste ( 61,3 ), e a menor, na região Sul ( 58,8 ).
A Mediana dos alunos de cursos Presenciais de Bacharelado, neste componente, foi 65,0, a mesma da região Centro-Oeste.
Nas regiões Norte ( 63,0 ), Nordeste ( 64,8 ) e Sul ( 64,0 ) a Mediana foi menor, enquanto na região Sudeste ( 65,5 ) foi maior.
A nota Máxima ( 98,5 ) foi a mesma na região Sudeste.
A nota Mínima ( 0,0 ) foi a mesma em todas as regiões do Brasil.
O Gráfico 6.6 representa a distribuição das notas nas questões discursivas do Componente de Formação Geral.
A moda desta distribuição ocorre no intervalo ( 60 ; 70 ], tanto para os cursos de Licenciatura quanto para os de Bacharelado.
Destaca -se também o intervalo [ 0 ; 10 ] com uma proporção acima de 5 \% do total de notas, para ambas as Habilitações, sendo que no intervalo [ 0 ; 10 ] incluem -se além da nota zero, a frequência de alunos que deixaram as respostas deste tipo de questão em branco.
Na sequência, os resultados verificados para cada uma das questões discursivas de Formação Geral serão apresentados, estabelecendo -se relações com os conteúdos abordados em cada uma delas.
Os comentários da Banca de docentes corretores, a respeito do observado na correção das respostas dos estudantes, suas impressões e conclusões serão apresentados junto à análise de cada questão.
Cumpre esclarecer que, tendo em vista que as questões discursivas de Formação Geral são padronizadas, ou seja, constam de todas as provas, os comentários da Banca são os mesmos para todas as carreiras acadêmicas, sendo direcionados a todos os estudantes que participaram do Enade/2017.
A seguir, serão analisados os desempenhos da Área de Ciências Biológicas ( Licenciatura/Bacharelado ) nas duas questões discursivas de Formação Geral do Enade/2017, comparando -se os resultados obtidos com comentários para cada questão. 6.3.1.1 Análise de Conteúdo da Questão Discursiva 1 do Componente de Formação Geral Os dados de Ciências Biológicas ( Licenciatura/Bacharelado ), obtidos a partir das respostas à questão 1, encontram -se na Tabela 6.14 e no Gráfico 6.7, desagregados por Habilitação e Modalidade de Ensino.
Nessa questão – de melhor desempenho dentre as duas de Formação Geral – os alunos, dos cursos de Licenciatura, tiveram Média 58,4.
A maior Média para a questão 1 entre os alunos de cursos de Licenciatura foi obtida na região Sudeste ( 59,9 ), e a menor, na região Norte ( 56,5 ).
Já os alunos dos cursos Presenciais de Licenciatura obtiveram uma Média mais alta, 58,9, e os de cursos a Distância, valor um pouco mais baixo, 55,9.
A maior Média para a questão 1 entre os alunos de cursos Presenciais de Licenciatura foi obtida na região Sul ( 60,6 ), e a menor, na região Norte ( 56,5 ).
Dentre os alunos de cursos a Distância de Licenciatura, a maior Média foi obtida na região Sudeste ( 58,9 ), e a menor, na região Centro-Oeste ( 51,4 ).
Quanto à variabilidade das notas, o Desvio padrão das notas foi 26,3 para os alunos de Educação Presencial e 26,5 para os de Educação a Distância.
Para os cursos Presenciais de Licenciatura, as medianas de quase todas as regiões e do Brasil, como um todo, foram iguais ( 65,0 ), com exceção das regiões Norte e Centro-Oeste com Mediana 60,0.
Já para os cursos a Distância, a Mediana do Brasil, como um todo, foi 60,0.
Nas regiões Norte, Nordeste, Centro-Oeste ( 55,0 em cada ) e Sul ( 57,5 ) as medianas foram mais baixas, enquanto na região Sudeste ( 65,0 ) a Mediana foi maior.
As notas máximas e mínimas da questão discursiva 1 foram para ambas as Modalidades de Ensino, sem exceção, as mesmas para todas as regiões do Brasil, respectivamente, 100,0 e 0,0.
Como já comentado, não foram observados cursos a Distância de Bacharelado.
Dentre os alunos de cursos Presenciais de Bacharelado, o desempenho nesta questão foi um pouco mais baixo do que na questão discursiva 2.
Os alunos dos cursos de Bacharelado obtiveram Média 60,8.
A maior Média para a questão 1 foi obtida na região Nordeste ( 62,1 ), e a menor, na região Sul ( 59,0 ).
Quanto à variabilidade das notas, o Desvio padrão das notas foi 26,1 para os alunos de Educação Presencial.
Para os cursos Presenciais de Bacharelado, as medianas, as notas máximas e as notas mínimas do Brasil e das regiões Norte e Sudeste foram as mesmas, 65,0, 100,00 e 0,0, respectivamente.
O Gráfico 6.7 mostra a distribuição das notas na questão discursiva 1 do Componente de Formação Geral.
Observa -se que a maior frequência corresponde ao intervalo ( 70 ; 80 ] para ambas as Habilitações.
Também é importante a proporção dos alunos que deixaram a resposta a essa questão em branco, com mais de 5 \% dos participantes de ambas as Habilitações. 6.3.1.2 Comentários sobre a correção de Conteúdo das respostas à Questão Discursiva 1 O enunciado era claro e indicava um tema de extrema relevância, qual seja, a incidência de doença grave integrante do rol das DST ( Doenças Sexualmente Transmissíveis ), que encaminha para a análise do direito fundamental à saúde, que tem assento constitucional ( CRFB, art. 6º. caput : São direitos sociais a educação, a saúde, a alimentação, o trabalho, a moradia, o transporte, o lazer, a segurança, a previdência social, a proteção à maternidade e à infância, a assistência aos desamparados, na forma desta Constituição ).
A questão focalizava, por meio de três textos, um tema social extremamente relevante, qual seja, a relação entre homem e mulher nas suas interações íntimas, no caso representadas pela vida sexual, e o reflexo das relações de poder veiculadas pela sociedade, que apontam, no mais das vezes, o predomínio do homem.
Tal posição de superioridade é justificada pelas convenções arraigadas que acarretam a educação dos homens e mulheres por padrões patriarcais, que redundam em posições machistas.
Um dos textos demandava uma reflexão sobre a realidade dos casais que convivem em relações estáveis ou fugazes e não têm o controle dos efeitos das relações sexuais, que abrangem a gravidez não planejada e, em decorrência do não cuidado, a inserção de doenças transmitidas sexualmente.
Esse panorama é refletido na pouca influência das mulheres sobre os homens quanto à proteção, normalmente representada pelo uso de preservativos nas relações intimas, e após a constatação da doença, no encaminhamento das pessoas para tratamento.
Como efeito colateral, está a transmissão de doenças para os fetos.
Outro texto trazia à tona o maior cuidado que as mulheres têm, em geral, com sua saúde, notadamente no momento em que se descobre grávida e que redunda na realização de exames obrigatórios de pré-natal, que podem facilmente detectar doenças e, caso encontradas, ter o seu tratamento encaminhado.
A par disso, destaca -se que a crise de recursos financeiros agravou a situação, que já recebe influxos sociais negativos, como a dificuldade do homem em utilizar preservativos, realizar exames e procurar tratamentos quando adoece.
Como muito divulgado pelas mídias, a crise financeira do país repercutiu diretamente na prestação de serviços de saúde em todas as esferas : distrital, federal, estadual e municipal.
Assim, agravou -se o problema da falta de medicamentos essenciais para o tratamento da sífilis e, certamente, de outras doenças.
Como se verifica, a temática indicada no enunciado, representada pelos textos, pertence à realidade conhecida amplamente pelos estudantes, pois também é divulgada por todos os meios de comunicação.
Trata -se de tema geral que, no entanto, normalmente, não se inclui nas grades curriculares de cursos superiores que não sejam da área de Saúde.
Apesar disso, a temática é de conhecimento geral e relacionada aos direitos humanos, esses sim, ponto de integração com as diretrizes curriculares.
A linguagem utilizada era clara, indicando a existência dos fatos vinculados a questões relevantes e à saúde pública.
No entanto, como os textos trazem uma ampla gama de possibilidades de interpretação, ocorreram algumas compreensões diversas das esperadas.
Além disso, ressalta -se que a amplitude dos textos disponibilizados no enunciado fornece elementos que poderiam ser simplesmente reproduzidos pelos concluintes nas suas respostas.
O enunciado da questão encaminha a resposta para dois aspectos.
O estudante deveria abordar conteúdo pertinente à vulnerabilidade das mulheres às DST e o papel social do homem em relação à prevenção dessas doenças.
Para responder a esses aspectos, um número expressivo de concluintes utilizou somente os elementos dos textos motivadores, sem nada acrescentar.
A resposta deveria conter também uma segunda parte, a indicação de duas ações voltadas para o público masculino, a serem adotadas no âmbito das políticas públicas de saúde ou de educação, para a redução do problema.
Verificou -se que o desempenho dos estudantes nessa questão foi de médio para bom.
Ao se analisarem as respostas, constata -se que muitos foram os que apresentaram os tópicos presentes no padrão de resposta, amplo o suficiente para albergar as respostas apresentadas pelos concluintes.
Tal resultado pode ser justificado, tendo em vista tratar -se de tema bastante difundido nas diversas mídias, por ter sido apresentado de forma ampla pelos textos motivadores presentes no enunciado.
Por outro lado, o debate no ambiente universitário sobre as relações de poder entre homens e mulheres tende a ser marcante pelo movimento de empoderamento feminino, liderado por organizações sociais com ampla atuação nos estabelecimentos de nível superior.
De forma recorrente, as relações patriarcais entre homens e mulheres foram vinculadas ao modelo machista vigente.
Essas condições, segundo os concluintes, explicam a incidência da doença e a situação de vulnerabilidade das mulheres, bem como a função social dos homens.
Outras respostas apontaram o poder público, nas áreas de saúde e/ou educação, como órgãos responsáveis pela solução do problema.
Para os que se fixaram na área da saúde, foram frequentes a indicação da necessidade de mais medicamentos, vacinas, educação sexual, médicos e postos de atendimentos.
Alguns, no entanto, apenas indicaram ações individuais para o homem, e outros, ações individuais para as mulheres.
Também figuraram respostas que apresentavam ações da igreja como solução, além de outras que faziam referências a campanhas no seio das famílias. 6.3.1.3 Análise de Conteúdo da Questão Discursiva 2 do Componente de Formação Geral A Tabela 6.15 apresenta as estatísticas básicas das notas dos alunos na questão discursiva 2 da área de Ciências Biológicas por Grande Região, segundo Habilitação e Modalidade de ensino.
Como já comentado, em particular, observa -se que o desempenho médio dos estudantes de cursos de Ciências Biológicas de Licenciatura, na questão discursiva 2, foi inferior ao obtido na questão discursiva 1 : médias, respectivamente 58,4 na questão 1 e 56,6 na questão 2.
Entre os cursos de Licenciatura, a região Sudeste foi aquela cuja Média nessa questão foi maior ( 59,0 ), e a região Nordeste foi a de menor Média ( 54,2 ).
Quando se considera também a Modalidade de Ensino, a Tabela 6.15 mostra que o desempenho médio dos estudantes de cursos Presenciais de Licenciatura na questão discursiva 2 ( média 57,4 ) foi também inferior ao obtido na questão discursiva 1 ( média 58,9 ).
A região Sul foi aquela cuja Média nessa questão foi maior ( 60,9 ), e a de menor Média foi a região Norte ( 54,6 ).
Quanto à variabilidade das notas, o Desvio padrão de todo o Brasil foi 31,7.
O maior desvio nessa questão foi obtido na região Nordeste ( 32,4 ), enquanto o menor foi obtido na região Sul ( 30,2 ).
As medianas 70,0, as notas máximas 100,0 e as notas mínimas 0,0 foram as mesmas em todas as regiões, sem exceção.
A média das notas dos estudantes de cursos a Distância de Licenciatura na questão discursiva 2 foi 52,3, também inferior à média obtida na questão discursiva 1 ( 55,9 ).
A região Sudeste foi aquela cuja Média nessa questão foi maior ( 55,5 ), e a região Norte foi a de menor Média ( 46,9 ).
O Desvio padrão de todo o Brasil foi 31,9, sendo que o maior desvio nessa questão foi obtido na região Nordeste ( 33,3 ), enquanto o menor foi obtido na região Norte ( 27,9 ).
A Mediana para esses alunos de cursos a Distância de Licenciatura foi 50,0, a mesma de quase todas as regiões, exceto a das regiões Norte ( 45,0 ) e Sudeste ( 70,0 ).
As notas máximas ( 100,0 ) e as notas mínimas ( 0,0 ) foram as mesmas em todas as regiões.
Como já comentado, não foram avaliados cursos a Distância de Bacharelado.
Dentre os alunos de cursos Presenciais de Bacharelado, o desempenho nessa questão foi um pouco melhor do que na questão discursiva 1 ( 60,8 ).
Os alunos dos cursos de Bacharelado, obtiveram Média 60,9.
A maior Média para a questão 2 foi obtida na região Sudeste ( 62,3 ), e a menor, na região Sul ( 58,9 ).
Quanto à variabilidade das notas, o Desvio padrão das notas foi 30,8 para os alunos de Educação Presencial.
Para os cursos Presenciais de Bacharelado, as medianas, as notas máximas e as notas mínimas do Brasil e de todas as regiões foram as mesmas : 70,0, 100,00 e 0,0 respectivamente.
O Gráfico 6.8 mostra a distribuição das notas na questão discursiva 2 do Componente de Formação Geral.
Observa -se que a maior frequência corresponde aos alunos que obtiveram nota dentro do intervalo ( 70 ; 80 ] para ambas as Habilitações.
Destacam -se, também, os alunos que deixaram a resposta dessa questão em branco, com proporção em torno de 10 \% para ambas as Habilitações. 6.3.1.4 Comentários sobre a correção de Conteúdo das respostas à Questão Discursiva 2 O enunciado era claro e abordava um tema de extrema relevância : a questão de gênero.
Trata -se de uma proposta polêmica, pois abrange a análise da inserção social de minoria vinculada à questão de gênero.
O enunciado solicitava que o concluinte discorresse sobre a importância do nome para as pessoas transgêneras e propusesse uma medida, no âmbito das políticas públicas, para facilitar o acesso dessas pessoas à cidadania.
A respeito da primeira parte, a importância do nome, esperava -se que o estudante abordasse as dificuldades que o uso do nome civil ( nome de registro, nome de batismo, nome da certidão de nascimento, por exemplo ) pode acarretar para as pessoas transgêneras, além de demonstrar conhecimento sobre a importância do nome social, direito já reconhecido através, principalmente, de decisões judiciais e atos administrativos, como decretos e portarias que dispõem sobre o uso do nome social.
Em atendimento à segunda parte do comando, o participante deveria propor uma política pública associada ao direito à cidadania das pessoas transgêneras.
Sobre esse aspecto, esperava -se a compreensão do significado de política pública, ou seja, propostas de âmbito governamental e não iniciativas de Organizações Sociais, igrejas ou grupos sociais de qualquer tipo.
O nascimento das pessoas, no Brasil, é objeto de inscrição em cartórios de Registro Civil das Pessoas Naturais, órgão responsável por estabelecer a certificação pública dos dados de todos os indivíduos que não somente nascem, mas casam, eventualmente se divorciam e vêm a óbito.
Daí o nome civil.
Para flexibilizar a correção, outras nomenclaturas foram consideradas, a saber : identificação de nascimento ; nome de batismo ; nome de registro, nome próprio, nome verdadeiro, que diferenciam o nome civil daquele que o indivíduo deseja utilizar para traduzir a sua opção pessoal de gênero, que se convencionou denominar nome social.
Quanto ao entendimento do termo transgênero, deve -se destacar que o movimento de mudança do gênero é contemporâneo e não está necessariamente vinculado à mudança de sexo.
Existem várias hipóteses em jogo.
Uma é a que envolve cirurgias de mudança de sexo ( a menina extirpa os seus seios e inclui órgão genital por intervenção cirúrgica ; o menino inclui seios e extirpa o seu órgão genital e inclui o órgão feminino também por intervenção cirúrgica ).
Há outro movimento que envolve apenas as vestimentas e o linguajar corporal, além dos relacionamentos ( as meninas assumem o trajar, o andar, os gostos dos meninos e mantêm relacionamentos afetivos com meninas ; os meninos assumem o trajar, o andar, os gostos das meninas e mantêm relacionamentos afetivos com meninos ).
Em relação à discussão da primeira solicitação do comando da questão, destaca -se que todos têm nome civil, que foi escolhido pelos seus pais ou responsáveis e, na falta deles, por autorização do Juiz de Direito, pautado no sexo detectado no nascimento.
Já o nome social é aquele que indica a preferência do indivíduo na sua identificação nas relações sociais e que destoa do seu gênero de registro civil.
Ressalte -se que o Brasil não possui ainda legislação específica sobre pessoas transgêneras.
Observa -se um relativo progresso no reconhecimento de direitos, através, principalmente, de decisões judiciais e atos administrativos, como decretos e portarias que dispõem sobre o uso do nome social.
Sem uma lei que defina os procedimentos da alteração dos documentos para pessoas transexuais, essa parcela da população LGBT ( lésbicas, gays, bissexuais, transexuais e transgêneros ) é obrigada a procurar na Justiça o reconhecimento de sua identidade, em processos que podem ser longos e que dependem do olhar de determinados profissionais como médicos, psicólogos, assistentes sociais, bem como do entendimento dos advogados, defensores públicos, juízes e do Judiciário como um todo.
Atualmente, tramita na Comissão de Direitos Humanos da Câmara dos Deputados o Projeto de Lei João Nery ( 5002/2013 ), dos deputados Jean Willys ( PSOL-RJ ) e Erika Kokay ( PT-DF ), que determina que o reconhecimento da identidade de gênero seja um direito do cidadão.
O projeto recebeu o nome do primeiro trans-homem operado no Brasil.
Grande parte dos respondentes considerou a incongruência entre a imagem corporal, representada pelo gênero assumido pelos trans, e o nome que apresenta nos seus documentos como causa de sérios constrangimentos e embaraços e as dificuldades para que esses sujeitos acessem direitos básicos como saúde, educação e empregabilidade numa sociedade heteronormativa e preconceituosa.
Muitos demonstraram reconhecer que a identidade e expressão de gênero, através do nome social, pressupõe que sejam respeitadas a dignidade da pessoa humana, a intimidade, a liberdade e o direito à cidadania plena.
No entanto, também foram encontradas, em quantidade significativas, respostas que demonstravam a existência de uma interpretação social conservadora que não admite a mudança de sexo e nem o nome social.
Existe um acentuado preconceito contra todos os movimentos que buscam identificação social diversa da civil.
Isso inclui os variados grupos : gays, travestis, transgêneros, etc.
Já em relação ao segundo tópico a ser abordado pelo concluinte, é notório que as pessoas trans encontram -se destituídas de direitos no campo das políticas públicas como saúde, habitação, emprego, educação, justiça, entre outros.
Nesse sentido, uma grande diversidade de medidas, no âmbito das políticas públicas, poderia ser proposta.
Ações voltadas para a capacitação de profissionais da educação, saúde, segurança pública entre outras áreas, para melhor atender às demandas e especificidades desse público alvo seriam necessárias.
Ainda em relação à educação, ações que contribuíssem para diminuir a evasão, pois se o ambiente escolar é difícil para uma pessoa transgênera, por conta das situações discriminatórias e preconceituosas vindas dos colegas de classe, e mesmo a discriminação oriunda das/dos professores e/ou de outros profissionais da área, não podem ser mais um complicador nesse processo.
Ainda no que tange ao assunto, percebe -se que ao serem excluídas dos bancos escolares, as oportunidades no mercado de trabalho formal praticamente inexistem.
Um outro aspecto que se evidencia no cotidiano desse segmento da sociedade diz respeito às situações de violência sofridas pelas pessoas transgêneras.
Nesse sentido, algumas propostas foram elencadas para minimizar o problema e combater a violência.
Políticas envolvendo uma melhor formação de funcionários das polícias civil e militar, por exemplo, somariam esforços na qualificação das abordagens policiais e dos atendimentos prestados pelas forças de segurança, pois o respeito à orientação sexual e à identidade e expressão de gênero é antes de tudo o respeito à dignidade, à cidadania e à própria democracia.
Nesse sentido, revela -se essencial que, ao sofrerem agressões, sejam elas físicas, psicológicas ou morais no espaço público ou domiciliar, as pessoas trans possam ser atendidas, com respeito inclusive ao nome social, nas delegacias, inclusive nas Delegacias Especiais de Atendimento à Mulher ( DEAM’s ), pois esse é o órgão institucional que possivelmente detém as qualificações necessárias para o atendimento ao gênero feminino ao qual as travestis e mulheres trans efetivamente pertencem.
Também ocorre uma latente inacessibilidade das pessoas trans à assistência jurídica.
Nesse sentido, é necessário implementar projetos de capacitação para profissionais do poder judiciário como um todo, de modo a prestarem a assistência jurídica gratuita na defesa intransigente dos direitos sociais, civis, trabalhistas e políticos dessa população.
Dessa maneira, a Justiça dar -se-ia de forma mais acessível e sensível para o segmento, que terá suas especificidades observadas a começar pelo devido respeito do direito ao uso do nome social.
Nesse âmbito, o mais comum foi propor a adoção de medidas punitivas para quem viola o direito à autodeterminação de gênero e a elaboração de leis que garantam a mudança do nome e assegurem outros direitos para as pessoas trans.
Propostas envolvendo a formação de profissionais, como as mencionadas aqui, e de conscientização da população foram as mais frequentes.
Muitas outras enfatizavam o direito ao uso do nome social respeitado e a garantida de alteração do registro civil, independentemente da cirurgia de redesignação sexual ou transgenitalização.
Muitos questionaram o fato de não bastar a autodeclaração no processo de alteração de nome, enfatizando que o nome social é uma questão de princípio de dignidade e cidadania.
Apesar de pouco mencionada nas respostas corrigidas, políticas envolvendo a promoção da saúde são muito necessárias.
Cabe aos gestores de saúde ações direcionadas, não somente à prevenção e tratamento do HIV/Aids, mas também políticas de saúde que contemplem as especificidades dessa população, tais como o Ministério da Saúde, 6.3.1.5 Análise de Língua Portuguesa das Questões Discursivas do Componente de Formação Geral Os dados obtidos a partir das respostas às questões discursivas do Componente de Formação Geral de Ciências Biológicas desagregados por Grande Região, segundo Habilitação e Modalidade de Ensino, no que tange à Língua Portuguesa, encontram -se na Tabela 6.16 e no Gráfico 6.9.
Nos cursos de Licenciatura em Ciências Biológicas, a Média desse aspecto foi 57,3, ligeiramente abaixo da obtida pelos alunos de Bacharelado ( 59,8 ).
Quando se considera a Modalidade de Ensino para cada Habilitação, os alunos de cursos Presenciais de Licenciatura obtiveram Média 57,7.
A Média em relação à Língua Portuguesa foi menor para os alunos de cursos a Distância ( 55,6 ).
Para os cursos de Bacharelado, como já comentado, não foram avaliados cursos a Distância de Bacharelado, e a Média foi 59,8 para Presenciais.
A Mediana das notas de Língua Portuguesa para os cursos de Licenciatura foi 62,5 para cursos Presenciais ( a mesma para o total ) e 60,0 para os cursos a Distância.
Para os cursos especialmente o SUS ( Sistema Único de Saúde ), incluir, de fato, o atendimento para pessoas trans para acompanhamento da terapia hormonal que muitas demandam, respeitando o direito de utilização do nome social.
De modo geral, são muitas as políticas públicas a sugerir e urge a necessidade de uma agenda mais compromissada por parte dos gestores do nosso país com essa comunidade que ainda é muito discriminada.
Logo, tais políticas devem contemplar suas necessidades básicas como o direito de acesso aos estudos, à capacitação para o mercado de trabalho, à profissionalização, ao acesso a bens e serviços de qualidade em saúde, habitação, segurança, cultura e assistência social e, certamente, o respeito ao uso do nome social, o que também contribui para essa inclusão.
Assim, as opções apresentadas no padrão de resposta funcionaram como exemplos de políticas públicas possíveis já que seria difícil apresentar uma lista suficientemente ampla para abranger a grande variedade de respostas encontradas.
O desempenho dos estudantes nessa questão foi médio.
As respostas, de modo geral, indicaram a necessidade de uma atualização por parte dos concluintes quanto a temas essenciais e, no caso da presente questão, a análise sobre aspectos fundamentais da questão de gênero.
Muitos, possivelmente motivados pelos textos apresentados no enunciado, restringiram -se a discorrer sobre preconceito de gênero ou mesmo sobre o preconceito de modo mais geral, sem responder ao que foi solicitado no comando da questão.
Outros apresentaram discursos religiosos ou ideológicos, também sem tratar do que foi solicitado. 6.3.1.5 Análise de Língua Portuguesa das Questões Discursivas do Componente de Formação Geral Os dados obtidos a partir das respostas às questões discursivas do Componente de Formação Geral de Ciências Biológicas desagregados por Grande Região, segundo Habilitação e Modalidade de Ensino, no que tange à Língua Portuguesa, encontram -se na Tabela 6.16 e no Gráfico 6.9.
Nos cursos de Licenciatura em Ciências Biológicas, a Média desse aspecto foi 57,3, ligeiramente abaixo da obtida pelos alunos de Bacharelado ( 59,8 ).
Quando se considera a Modalidade de Ensino para cada Habilitação, os alunos de cursos Presenciais de Licenciatura obtiveram Média 57,7.
A Média em relação à Língua Portuguesa foi menor para os alunos de cursos a Distância ( 55,6 ).
Para os cursos de Bacharelado, como já comentado, não foram avaliados cursos a Distância de Bacharelado, e a Média foi 59,8 para Presenciais.
A Mediana das notas de Língua Portuguesa para os cursos de Licenciatura foi 62,5 para cursos Presenciais ( a mesma para o total ) e 60,0 para os cursos a Distância.
Para os cursos Presenciais de Bacharelado, foi 65,0.
A nota máxima não chegou a 100,0 para nenhuma combinação de Habilitação e modalidade.
No Gráfico 6.9, mostra -se a distribuição das notas de Língua Portuguesa do Componente de Formação Geral.
Observa -se que a maior frequência ( em torno de 30 \% ) corresponde à dos alunos que obtiveram nota no intervalo ( 60 ; 70 ] em ambas as Habilitações.
Deve -se destacar, também, os alunos que deixaram as respostas de ambas as questões em branco representando em torno de 5 \% do total, caracterizando -se como máximo local. 6.3.1.6 Comentários sobre a correção das respostas de Formação Geral com respeito à Língua Portuguesa Os enunciados apresentados em 2017 não explicitaram a exigência de elaboração de um texto “ dissertativo ”, como nos anos anteriores.
O participante deveria fazer uma exposição de seus conhecimentos sobre cada um dos assuntos e estruturar seus textos de acordo com as características do registro formal adequado à situação comunicativa – avaliação de conhecimentos.
Essa configuração determina exigências quanto aos seguintes aspectos : adequação da seleção vocabular, desenvolvimento do conteúdo, organização lógica das ideias, estruturação sintática dos períodos, utilização de procedimentos de encadeamento textual e de referenciação, obediência às exigências morfossintáticas próprias da modalidade escrita da norma-padrão, respeito às regras ortográficas e às regras de acentuação gráfica.
O padrão de resposta utilizado na avaliação das questões 1 e 2 considerou os aspectos relevantes ao bom desempenho linguístico como competências distintas, de modo a permitir um mapeamento detalhado do domínio dos recursos disponíveis na Língua Portuguesa para a comunicação escrita formal.
Com base nesse objetivo, foram avaliados os seguintes aspectos : ( a ) estruturação textual condizente com o gênero solicitado e o modo de organização textual expositivo adequado ao gênero – essa competência envolve : • estruturação sintática condizente com o padrão da modalidade escrita formal da língua portuguesa de modo a garantir a clareza necessária ; • distribuição do conteúdo do texto em parágrafos, de modo a garantir a sua organização temática ; • utilização de operadores discursivos que contribuam para a progressão temática do texto, estabelecendo relações lógicas entre as ideias apresentadas, tanto do ponto de vista intrafrasal, como do interfrasal ; • utilização de procedimentos de referenciação lexical e pronominal que permitam a retomada de referentes textuais ; • utilização de sinais de pontuação que contribuam para a organização lógica da frase e do texto ; • inteligibilidade relacionada ao atendimento das exigências de estruturação textual.
Espera -se, portanto, que o participante recorra a procedimentos linguístico-discursivos para organizar seu texto, permitindo o encadeamento lógico entre suas partes de forma a garantir a progressão e a coerência textuais.
Isso significa que os seguintes procedimentos foram considerados inadequados, de acordo com o padrão de resposta proposto : • elaboração de frases fragmentadas que comprometam a estrutura lógico-gramatical do texto ; • sequência justaposta de ideias sem encaixamentos sintáticos, reproduzindo hábitos da oralidade ; • elaboração de frase com apenas oração subordinada, sem oração principal ; • emprego equivocado do conector ( preposição, conjunção, pronome relativo, alguns advérbios e locuções adverbiais ) comprometendo a expressão da relação lógica entre duas ideias, com prejuízo da clareza do texto ; • emprego do pronome relativo sem a preposição, quando obrigatória ; • repetição ou substituição inadequada de palavras sem o emprego dos recursos oferecidos pela língua ( pronome, advérbio, artigo, sinônimo ) ; • emprego inadequado dos pronomes relativos “ cujo(a) ” e “ onde ” ; • utilização inadequada dos sinais de pontuação, comprometendo a clareza textual ; • inteligibilidade relacionada ao atendimento das exigências de estruturação textual. ( b ) respeito às convenções ortográficas da norma-padrão da Língua Portuguesa – essa competência envolve o domínio das regras de acentuação gráfica e da grafia padrão das palavras ( com ausência de abreviaturas próprias da linguagem da internet ), de acordo com as convenções estabelecidas pela legislação em vigor e consubstanciadas no Vocabulário Ortográfico da Língua Portuguesa, editado pela Academia Brasileira de Letras. ( Vale notar que, nessa edição de 2017, não foi aceita a legislação anterior, no caso das regras relativas ao uso do hífen e da acentuação gráfica ).
Espera -se que o participante :  grafe corretamente as palavras ;  respeite as regras de acentuação gráfica ;  empregue maiúsculas em início de frase, em nomes próprios de pessoas, lugares ou instituições ;  grafe as siglas com letras maiúsculas ou apenas com a primeira letra maiúscula quando formarem uma palavra ;  evite abreviações como p/, vc, tb, pq, tá, né, usadas muitas vezes em escrita informal e na internet ;  obedeça às regras de separação de sílabas no final da linha. ( c ) domínio dos diferentes aspectos morfossintáticos próprios da modalidade escrita formal da norma-padrão da Língua Portuguesa – essa competência envolve : a concordância nominal, a concordância verbal, a regência nominal, a regência verbal, a flexão nominal, a flexão verbal, a correlação entre os tempos verbais, a colocação pronominal e a utilização de sinais de pontuação que contribuam para a organização lógica da frase e do texto.
Espera -se que o participante :  flexione o verbo para estabelecer concordância de número com o sujeito da frase ;  flexione o artigo, o adjetivo e o pronome para concordar em número e em gênero com o substantivo a que se referem ;  observe a regência nominal e a verbal, utilizando a preposição adequada depois de um substantivo, um verbo ou um adjetivo ;  empregue adequadamente o acento grave indicador da crase entre uma preposição e um artigo ( a+a ) ;  obedeça às regras de colocação pronominal ( próclise e ênclise ), distintas dos hábitos da oralidade ou da escrita informal ( exigência de próclise com termo atrator, não exigência de que o pronome oblíquo se ligue ao verbo auxiliar por meio de hífen ) ;  flexione adequadamente verbos, substantivos, adjetivos e pronomes no que diz respeito à expressão das categorias gramaticais ;  flexione os verbos para expressar a correlação de modo e tempo nas estruturas subordinadas.
Com base nesses critérios, foram considerados como desvios de caráter morfossintático, e não como desvios ortográficos, as alterações que envolvem mudança de classe gramatical ou de forma flexional do verbo :  eliminação da marca de infinitivo ( -r- ) e substituição por acento agudo ou ausência total de marca do infinitivo ;  confusão entre “ ão ” e “ am ” nas formas verbais ;  confusão entre “ há ” e “ a ” ;  uso de hífen para separar pronome átono – tanto uso indevido quanto omissão ( exemplo : “ esperasse ”, em lugar de “ espera -se ” ; “ falar -mos ”, no lugar de “ falarmos ” ;  “ esta ” ( no lugar de “ está ” ) ; “ mais ” ( no lugar de “ mas ” ) ; “ e ” ( no lugar de “ é ” ).  verbos “ ter ” e “ vir ” que, na terceira pessoa do plural, não apresentarem o acento circunflexo, serão considerados como desvio de concordância, em aspectos morfossintáticos. ( d ) seleção vocabular adequada à modalidade escrita formal da Língua Portuguesa, exigida pela situação comunicativa – essa competência envolve a precisão na seleção/utilização do vocabulário relacionado à temática solicitada pela questão ; a ausência de marcas da oralidade, como termos de sentido muito genérico ( “ coisa ”, “ negócio ”, “ você ” ) e termos de registros mais informais ( como gírias, jargões, frases feitas, ditados populares, termos regionais ).
Assim, espera -se que o participante respeite a adequação vocabular não usando gírias ou expressões coloquiais, evite repetição desnecessária de palavras e utilize um vocabulário mais formal, como solicitado por um texto dissertativo.
Observações : - A inteligibilidade relacionada ao atendimento das exigências de estruturação textual foi avaliada na segunda competência ( relativa aos aspectos textuais ). - Os problemas de coerência textual provocados por uso indevido do vocabulário foram avaliados na terceira competência ( relativa aos aspectos morfossintáticos e vocabulares ). - Não foram considerados como desvios morfossintáticos os problemas de caligrafia ( -a/ -o, -s, -r ). - Não foram considerados os textos grafados integralmente em caixa alta, já que não seria possível distinguir alguma marcação especial para as letras em início de frase. - Cada desvio foi considerado como uma ocorrência, mesmo que dois desvios fossem relativos ao mesmo aspecto linguístico.
A grade de avaliação do desempenho linguístico considerou, portanto, três grandes grupos de competências, segundo os aspectos explicitados anteriormente : a ) domínio das convenções ortográficas : grafia de vogais e consoantes, uso de maiúsculas e minúsculas, emprego do hífen e acentuação gráfica ; b ) domínio dos procedimentos de estruturação textual do ponto de vista microestrutural : organização interna dos períodos, emprego de conectores para a articulação lógica entre os períodos e entre os parágrafos, emprego de marcas de referenciação lexical e pronominal ; utilização dos sinais de pontuação que contribuem para a organização lógica da frase ; c ) domínio das regras de caráter morfossintático estabelecidas como modelares do ponto de vista da modalidade escrita formal da norma-padrão da Língua Portuguesa : concordância nominal e verbal, regência nominal e verbal, colocação pronominal, flexão nominal e verbal, correlação entre tempos e modos verbais, ausência de marcas de oralidade.
A seleção vocabular adequada à modalidade escrita formal da Língua Portuguesa foi incorporada a essa última competência, tendo em vista a intersecção entre as duas do ponto de vista das exigências do registro formal da modalidade escrita da norma-padrão.
Apreciação geral do desempenho dos estudantes : A correção revelou desempenhos distintos dos participantes, marcados pela falta de repertório cultural da maioria dos participantes : - a questão 1 propiciou a oportunidade de maior desenvolvimento do tema solicitado, gerando consequências na estruturação textual, já que os textos foram mais longos, mais elaborados e fluentes.
Entretanto, em virtude da diversidade de abordagens dos textos motivadores ( a falta de penicilina, a recusa no uso da camisinha pelos homens, a maior frequência das mulheres na realização de exames ginecológicos, a transmissão das DST aos bebês durante a gestação ), observaram -se respostas que se configuraram como verdadeiras paráfrases, com comprometimento do caráter autoral desejável em uma questão discursiva de Formação Geral.
Alguns participantes selecionaram um trecho de cada texto motivador para compor suas respostas, outros copiaram trechos inteiros. - a questão 2 teve um comportamento distinto, devido ao caráter polêmico que o caracteriza.
Assim, os textos são mais curtos, com vocabulário repetitivo e preso aos textos motivadores.
Além disso, observou -se grande ocorrência de protesto em função do tema, ora por motivos religiosos ora por preconceito de gênero.
Quanto aos aspectos linguísticos analisados durante esta avaliação, os resultados observados estão descritos nos parágrafos seguintes.
Aspectos ortográficos : O desempenho dos participantes revelou uma diferença muito grande nos dois aspectos analisados nesta competência : baixo índice de desvios da grafia padrão e grande índice de desvios de acentuação.
Em vários casos, ocorre ausência completa de acentuação gráfica.
Os resultados revelam que a tendência dominante entre os universitários brasileiros é a eliminação da acentuação gráfica, talvez motivada pelos hábitos relacionados às redes sociais e pela ausência de esclarecimento dos meios de comunicação, das autoridades e das escolas sobre as decisões do Acordo Ortográfico da Língua Portuguesa de 1990.
Os casos mais sistemáticos de eliminação do acento indicador da sílaba tônica são : - palavras proparoxítonas ( “ sifilis ”, “ proximos ”, “ politicas ”, “ publicas ”, “ transgeneros ” ) ; - palavras paroxítonas terminadas em ditongo crescente ( “ necessario ”, “ noticiarios ”, “ individuo ”, “ dependencia ”, “ varios ”, “ propria ”, “ transmissiveis ”, “ ocorrencia ” ) ; - palavras paroxítonas com hiato ( “ saude ” ) - palavras oxítonas ( “ ninguem ”, “ esta ”, “ ate ”, “ tambem ”, “ prevencao ” ).
Por outro lado, destaca -se o uso indevido do acento gráfico em palavras como “ genêro ”, “ melâncolia ”, “ prevênção ”, “ intervênção ”, por exemplo.
Há, também, uma tendência a acentuar paroxítonos como se fossem oxítonos : “ transmissivéis ”. elaborados e fluentes.
Entretanto, em virtude da diversidade de abordagens dos textos motivadores ( a falta de penicilina, a recusa no uso da camisinha pelos homens, a maior frequência das mulheres na realização de exames ginecológicos, a transmissão das DST aos bebês durante a gestação ), observaram -se respostas que se configuraram como verdadeiras paráfrases, com comprometimento do caráter autoral desejável em uma questão discursiva de Formação Geral.
Alguns participantes selecionaram um trecho de cada texto motivador para compor suas respostas, outros copiaram trechos inteiros. - a questão 2 teve um comportamento distinto, devido ao caráter polêmico que o caracteriza.
Assim, os textos são mais curtos, com vocabulário repetitivo e preso aos textos motivadores.
Além disso, observou -se grande ocorrência de protesto em função do tema, ora por motivos religiosos ora por preconceito de gênero.
Quanto aos aspectos linguísticos analisados durante esta avaliação, os resultados observados estão descritos nos parágrafos seguintes.
Aspectos ortográficos : O desempenho dos participantes revelou uma diferença muito grande nos dois aspectos analisados nesta competência : baixo índice de desvios da grafia padrão e grande índice de desvios de acentuação.
Em vários casos, ocorre ausência completa de acentuação gráfica.
Os resultados revelam que a tendência dominante entre os universitários brasileiros é a eliminação da acentuação gráfica, talvez motivada pelos hábitos relacionados às redes sociais e pela ausência de esclarecimento dos meios de comunicação, das autoridades e das escolas sobre as decisões do Acordo Ortográfico da Língua Portuguesa de 1990.
Os casos mais sistemáticos de eliminação do acento indicador da sílaba tônica são : • palavras proparoxítonas ( “ sifilis ”, “ proximos ”, “ politicas ”, “ publicas ”, “ transgeneros ” ) ; • palavras paroxítonas terminadas em ditongo crescente ( “ necessario ”, “ noticiarios ”, “ individuo ”, “ dependencia ”, “ varios ”, “ propria ”, “ transmissiveis ”, “ ocorrencia ” ) ; • palavras paroxítonas com hiato ( “ saude ” ) • palavras oxítonas ( “ ninguem ”, “ esta ”, “ ate ”, “ tambem ”, “ prevencao ” ).
Por outro lado, destaca -se o uso indevido do acento gráfico em palavras como “ genêro ”, “ melâncolia ”, “ prevênção ”, “ intervênção ”, por exemplo.
Há, também, uma tendência a acentuar paroxítonos como se fossem oxítonos : “ transmissivéis ”.
Quanto ao domínio das convenções relativas à grafia das palavras, observam -se desvios recorrentes, como a hipercorreção pela escolha de “ e ” no lugar de “ i ”, por influência de hábitos da oralidade ( “ descriminação ” no lugar de “ discriminação ”, “ entervenção ” no lugar de “ intervenção ” ) ; desvios de grafia relacionados à variação diastrática podem ser observados em “ estrupo ”, “ subjulgadas ”, “ precoseito ”, “ soubre ”, “ vecendo ”, “ indesencia ”, “ apolojia ”, “ fulga ”, “ dereitos ”, “ sifelis ”.
Observam -se, também, muitos casos de inadequação no uso da maiúscula para destacar determinadas palavras-chave do texto, como “ Transgêneros ”, “ Brasileiros ”, “ Homens ”.
Destaque -se, também, o grande número de participantes que grafam os textos inteiramente em caixa alta.
Vale observar, também, que, ao contrário do que se esperava, não apareceram abreviaturas próprias do “ internetês ”, relacionadas ao uso de redes sociais e e-mails.
Aspectos textuais : Esta competência é a que se revela como a mais problemática entre os participantes, porque são muitos os problemas observados, desvios acumulados durante toda a formação do estudante e que não se resolvem com um estudo autodidata, como acontece com regras ortográficas ou morfossintáticas.
São eles : sequência justaposta de ideias sem encaixamentos sintáticos ; redução drástica de estruturas subordinadas, ao lado do aumento na frequência de estruturas coordenadas e absolutas ; redução no uso de conectores para expressar relações lógicas essenciais à construção do texto, substituídas pela exigência de inferência por parte do interlocutor para suprir a sua ausência ; emprego equivocado de operadores que não estabelecem relações lógicas coerentes entre ideias do texto ; emprego inadequado do pronome relativo ( com omissão da preposição ou a utilização de pronome inadequado, como “ onde ” ) ; repetição exaustiva de termos sem a utilização de procedimentos mais sofisticados de substituição ( hiperonímias, hiponímias, nominalizações, expressões metafóricas ) ; frases fragmentadas que comprometem a estrutura lógico-gramatical ; frases formadas apenas por oração subordinada, sem oração principal.
Um importante aspecto a destacar é o baixo desempenho de uma parte dos participantes em relação à estrutura formal do texto produzido, o que é preocupante ao se levar em conta que são graduandos em fase final de formação.
São frequentes os casos de desvios de estruturação frasal, com uso inadequado ou ausência de conectivos entre parágrafos e entre frases.
Em uma parte dos textos, falta textualidade e domínio do registro padrão da língua.
Na verdade, observam -se relações linguísticas quase agramaticais, como as estabelecidas pela sequência de gerúndios sem o apoio de um ponto de partida para a organização das informações gramaticais e semânticas.
Observou -se que uma parte dos participantes não distribuiu as ideias em parágrafos, talvez devido ao pequeno número de linhas disponibilizadas para a resposta da questão ou, quem sabe, pela suposição de que não seria necessária essa divisão por não se tratar de um texto no modelo de uma redação dissertativo-argumentativa, como solicitado nos vestibulares.
Em função do encaminhamento dos enunciados das duas questões, que solicitaram uma análise do problema e encaminhamentos de políticas públicas, houve uma grande tendência, também, de construção de dois parágrafos desconexos, sem utilização de elementos coesivos adequados à progressão textual.
Quanto à utilização dos mecanismos de referenciação, deve -se destacar a ocorrência de repetições de palavras ou expressões sem a utilização de termos sinônimos ou pronomes, como seria adequado.
Quanto à utilização dos sinais de pontuação, observou -se uma grande precariedade nos textos analisados.
É muito frequente a ocorrência de parágrafos sem marca interna de pontuação para separar os períodos.
Vale observar que não foi penalizada a ausência de vírgula para destacar locuções ou adjuntos adverbiais de pequena extensão deslocados de posição na frase, por ser um uso opcional.
São os seguintes os tipos de problemas encontrados : a ) vírgula : utilização de vírgula para separar o sujeito e o predicado ; ocorrência de apenas uma das vírgulas para separar uma palavra, uma expressão ou uma oração encaixada ; uso de vírgula no lugar do ponto para separar ideias que constituem períodos distintos ; ausência de vírgula para separar elementos de uma enumeração ; ausência de vírgula para separar oração adjetiva explicativa ou utilização inadequada para separar oração adjetiva restritiva ; b ) ponto e vírgula : utilização do ponto e vírgula no lugar de vírgula ; c ) ponto final : ausência de ponto final para separar períodos.
Aspectos morfossintáticos e vocabulares : Em relação à regência, o desvio mais frequente é a falta do sinal indicativo da crase – isso revela que o usuário não tem consciência de que, sob a forma do termo “ a ”, existe a presença de uma contração entre a preposição “ a ” ( exigida pela regência do termo anterior ) e o artigo definido “ a ”.
Um desvio de regência significativo, nos últimos anos, é a utilização inadequada de uma preposição ou sua ausência após o verbo ou o nome ( substantivo ou adjetivo ).
Outro problema relacionado à regência verbal e à nominal, encontrado frequentemente nas questões, foi a ausência de preposição antes de pronome relativo, processo generalizado na modalidade oral da língua, em situações de registro informal.
Apesar da possibilidade de que essa alteração de regência se generalize no padrão escrito da Língua Portuguesa, como já está ocorrendo até em textos jornalísticos, o não emprego da preposição foi considerado inadequado neste processo de avaliação.
A concordância verbal e a concordância nominal apresentam alguns desvios muito frequentes.
Quanto à concordância de número, observou -se ausência de marca ( com sujeito anteposto ou posposto ) ou uso indevido ( uso inadequado da marca de plural comandado pelo núcleo plural da locução adjetiva, apesar de o substantivo que funciona como núcleo do sintagma nominal estar no singular ).
Uma ocorrência que se destacou foi a ausência de acento circunflexo na forma plural do presente do indicativo do verbo “ ter ”, que foi considerada como um desvio na concordância verbal e não na acentuação gráfica.
Quanto à concordância de gênero, vários casos foram observados, normalmente no âmbito de sintagmas nominais longos, em que o adjetivo está afastado do substantivo.
Deve -se destacar o aparecimento da marca de plural em verbos ou adjetivos comandados por adjunto adnominal plural, apesar de serem relacionados a núcleos substantivos no singular, evidenciando um processo de hipercorreção ( exemplo : “ O nome social dos trangêneros podem ser estabelecidos por uma legislação específica ” ).
Quanto à questão da colocação pronominal, foram poucos os desvios observados.
Concluiu -se que, no registro escrito formal, a maioria dos participantes já incorporou regras como a não introdução da frase por um pronome oblíquo e a próclise na presença de um termo atrator.
Não se adotou, entretanto, o padrão excessivamente formal descrito pelas gramáticas normativas em relação à posição do pronome oblíquo em locuções verbais, já que esse uso está muito distante da prática cotidiana, até em textos mais formais.
Quanto aos aspectos vocabulares, alguns tipos de inadequação foram observados : expressões da oralidade ; seleção vocabular incompatível com o contexto, gerando falta de inteligibilidade ; falta de domínio de vocabulário mais abstrato e de maior complexidade, essencial ao desenvolvimento do texto de base dissertativa.
O principal aspecto observado foi a excessiva repetição de certas palavras, revelando limitação de repertório vocabular.
O termo “ pessoa ”, por exemplo, chega a ser repetido até 7 ou 8 vezes em um mesmo texto.
Várias marcas de oralidade foram identificadas, embora não com alta frequência : uso do pronome relativo “ onde ” como relativo universal, falta de artigo definido antes de substantivo, falta de artigo definido antes de substantivo, reduções como “ tá ”, “ pra ”, “ pro ”, “ prum ”, expressões informais, eliminação de preposições.
Em função do tema solicitado na questão 1, a sigla DST foi grafada de diferentes maneiras : DST´s, DSTs ou DSTS.
Além disso, houve flutuação de gênero gramatical nos sintagmas que envolviam a sigla : “ os diferentes DSTs ” ou “ as diferentes DSTs ”.
Em função do tema solicitado na questão 2, houve flutuação de gênero gramatical nos sintagmas que envolviam a sua palavra-chave : “ os trans ” ou “ as trans ” ; “ os transgêneros ” ou “ as transgêneras ”.
Todas essas formas foram aceitas, por considerarmos que sua utilização é muito recente e os participantes ainda não sistematizaram essas formas em seu uso da língua. 6.3.2 Componente de Conhecimento Específico As estatísticas básicas das notas referentes às questões discursivas do Componente de Conhecimento Específico de Ciências Biológicas Licenciatura por Grande Região, segundo Modalidade de Ensino, estão disponibilizadas na Tabela 6.17a.
As informações equivalentes para o Bacharelado estão na Tabela 6.17b.
Observa -se que as Médias foram bem mais baixas do que para as questões discursivas do Componente de Formação Geral.
Enquanto no Componente de Formação Geral as Médias para estudantes de Ciências Biológicas de Licenciatura foram, respectivamente para os de Ensino Presencial e a Distância, 58,0 e 54,3 ( 57,4 para o total ), na parte de Conhecimento Específico as Médias foram 35,4 e 32,6 ( 34,9 para o total ).
Como já comentado, não foram observados cursos a Distância de Bacharelado.
Para os cursos Presenciais de Bacharelado, as médias equivalentes foram 60,6 para as discursivas de Formação Geral e 23,2 para as de Conhecimento Específico.
Para os cursos de Licenciatura, a nota Máxima foi maior para os cursos Presenciais ( 100,0 ) do que para os a Distância ( 96,7 ).
Para os cursos Presenciais de Bacharelado, a nota Máxima foi 100,0.
As notas mínimas foram zero para todas as combinações de Habilitação e Modalidade de Ensino.
O Gráfico 6.10 representa a distribuição das notas nas questões discursivas no Componente de Conhecimento Específico.
A moda desta distribuição ocorre no intervalo [ 0 ; 10 ], com cerca de 20 \% do total de participantes de Licenciatura e mais do que 30 \% do total de participantes de Bacharelado.
Observa -se que nesses intervalos estão computadas as ocorrências de respostas em branco para as três questões discursivas de Conhecimentos Específicos.
A análise de cada uma dessas questões será feita a seguir : 6.3.2.1 Análise da Questão Discursiva 3 do Componente de Conhecimento Específico A questão 3 foi a mesma para os estudantes de Bacharelado e de Licenciatura.
Os resultados aferidos encontram -se descritos nas Tabelas 6.18a e 6.18b, respectivamente para os cursos de Licenciatura e de Bacharelado.
As Médias dos estudantes de cursos de Licenciatura foram 6,2 e 4,2, respectivamente para os cursos Presenciais e os a Distância ( 5,8 para o total ), ao passo que a Média dos estudantes de cursos Presenciais de Bacharelado foi 9,8.
Essa foi a questão de desempenho mais baixo dentre as três discursivas de Conhecimentos Específicos para ambas as modalidades.
Para todas as combinações de Habilitação e modalidade, a nota Máxima foi 100,0, a Mínima foi 0,0, e a Mediana foi 0,0, as mesmas em quase todas as regiões do Brasil.
Isso indica que, independentemente da Habilitação e da Modalidade de Ensino, pelo menos metade dos estudantes de Ciências Biológicas de todo o Brasil tirou nota zero nessa questão.
O Gráfico 6.11 mostra a distribuição das notas na questão discursiva 3, do Componente de Conhecimento Específico da área de Ciências Biológicas ( Licenciatura/Bacharelado ).
A moda dessa distribuição é o intervalo das notas zero, com cerca de 45 \% do total de participantes em ambas as Habilitações.
Destaca -se o fato de que cerca de 85 \% do total de participantes de Licenciatura e de cerca de 75 \% do total de participantes de Bacharelado ou deixaram a respostas dessa questão em branco ou tiraram nota zero. 6.3.2.2 Comentários sobre as respostas à Questão Discursiva 3 A Questão 3 foi de fácil leitura, objetiva, sendo o seu conteúdo abordado no primeiro ano do ciclo básico, na maioria das IES.
O enunciado estava claro, com suficiência de informações para a resposta correta do estudante.
O resultado demonstra que faltou conhecimento ou houve esquecimento do conteúdo por parte dos graduandos do Bacharelado e da Licenciatura.
Não houve interpretações diferentes do padrão adotado.
O conteúdo abordado, segundo as Diretrizes Curriculares do MEC para a área de Ciências Biológicas, encontra -se em “ Biologia celular, molecular e evolução : visão ampla da organização e interações biológicas, construída a partir do estudo da estrutura molecular e celular, função e mecanismos fisiológicos da regulação em modelos eucariontes, procariontes e de partículas virais, fundamentados pela informação bioquímica, biofísica, genética e imunológica.
Compreensão dos mecanismos de transmissão da informação genética, em nível molecular, celular e evolutivo ” e, portanto, bastante pertinente e importante para os graduandos em Ciências Biológicas ( Bacharelado e Licenciatura ).
O nível de profundidade exigido foi o básico para essa área, imprescindível para a formação em Ciências Biológicas, mas foi a questão de maior dificuldade para os estudantes.
Por essa razão, foi considerada difícil em relação às outras questões de Conhecimentos Específicos.
Por não haver como tangenciar a resposta, a questão não dava margem para elaboração de respostas sem bom conhecimento do conteúdo exigido.
A abordagem usada para responder a essa questão objetiva só poderia ser uma.
Portanto, o número de questões em branco foi elevado.
Devido ao tipo de questão, o padrão de resposta era pouco flexível.
A flexibilização prevista no padrão dizia respeito considerar glicólise anaeróbica ou fermentação.
No caso de fermentação, muitos não ressaltaram a ausência de oxigênio ou condições anaeróbicas.
Tal resposta, ainda que incompleta, poderia ser pontuada.
A opção mais frequente foi deixar em branco.
A segunda mais frequente foi responder, apenas, ao item ‘ a ’ de forma totalmente absurda, tal como mencionar a fotossíntese como processo e ausência de mitocôndrias em plantas, que apenas possuem cloroplastos e não fazem respiração.
Hoje em dia, em particular nas IES Públicas, os estudantes de Bacharelado têm buscado a Licenciatura.
Assim, a antiga separação de conhecimentos entre os dois cursos tem diminuído muito, com os bacharéis pedindo matrícula para realizarem a Licenciatura.
No entanto, os estudantes de Bacharelado respondem mais e melhor às questões de conteúdo do que os da Licenciatura.
Os estudantes não foram bem nessa questão.
O desempenho geral tanto no Bacharelado como na Licenciatura foi bastante insatisfatório e não houve muita diversidade nas respostas apresentadas.
No entanto, foram observadas respostas absurdas, impressionantes para um graduando de Ciências Biológicas, como"as plantas não respiram","se não tem mitocôndria, tem cloroplasto","fotossíntese", etc.
Dentre as respostas corrigidas, desconsiderando as que estavam em branco, tanto do Bacharelado quanto da Licenciatura, as que obtiveram notas de zero a 30 foram a maioria ( 81 \% e 84 \%, respectivamente ).
Fazem parte desse grupo aqueles que não sabiam distinguir os processos vitais ( fermentação, fotossíntese, respiração aeróbica ) e desconheciam as funções das organelas e noções básicas de bioquímica celular.
Mais uma vez, desconsiderando as respostas em branco, as consideradas medianas ( notas de 35 a 70 ) foram, aquelas que responderam ao item ‘ a ’, demonstrando conhecer o nome do processo anaeróbico : glicólise anaeróbica ou fermentação.
De um modo geral, os estudantes \textbf{localizaram} o processo como sendo o citoplasma.
Os resultados foram similares entre os graduandos do Bacharelado e da Licenciatura.
As respostas boas ( notas de 75 a 100 ) foram muito poucas, tanto no Bacharelado quanto na Licenciatura ( desconsiderando as em branco, 5 \% e 3 \%, respectivamente ), tendo sido atribuídas aos que responderam aos dois subitens de modo claro, manifestando conhecimento do conteúdo em bioquímica e biologia celular.
Algumas ponderações devem ser feitas.
A Questão 3 era de conteúdo geral e básico.
Os alunos não pareceram sedimentar correta e satisfatoriamente os conteúdos abordados nos primeiros anos da graduação, e que são fundamentais para a sua formação.
São os primeiros tijolos de uma construção e, infelizmente, estão parecem estar mal estruturados na formação doa participantes do Enade/2017. 6.3.2.3 Análise da Questão Discursiva 4 do Componente de Conhecimento Específico As questões 4, presentes nas provas de Bacharelado e de Licenciatura, eram iguais.
As Tabelas 6.19a e 6.19b contêm as informações relativas à questão 4 do conjunto de questões do Componente de Conhecimento Específico.
O desempenho dos estudantes nessas questões foi superior ao da questão 3 e inferior ao da questão 5, para ambas as Habilitações.
As Médias para os alunos de Licenciatura ( ver Tabela 6.19a ) foram, respectivamente para os de cursos Presenciais e a Distância, 42,6 e 38,9 ( 42,0 para o total ).
A nota Máxima ( 100,0 ) foi atingida por, pelo menos, um concluinte de todas as regiões para ambas a Modalidades de Ensino.
A Mediana do Brasil, como um todo, foi 40,0 para ambas as Modalidades de Ensino.
Nos cursos Presenciais, as medianas foram : 0,0 na região Norte, 20,0 na região Nordeste, 60,0 nas regiões Sudeste e Sul e 40,0 na região Centro-Oeste.
Já nos cursos a Distância, as medianas foram : 0,0 nas regiões Norte, Nordeste e Centro-Oeste, e 40,0 nas regiões Sudeste e Sul.
Destaca -se o fato de que nas regiões Norte, para ambas as modalidades, e nas regiões Nordeste e Centro-Oeste, para os cursos a Distância, pelo menos a metade dos estudantes tiraram nota zero nessa questão.
A nota mínima foi zero para todas as combinações de Modalidade de Ensino e região.
As Médias para os alunos de cursos Presenciais de Bacharelado ( ver Tabela 6.19b ) foi 28,1.
A nota Máxima ( 100,0 ) foi atingida por, pelo menos, um concluinte de todas as regiões.
A Mediana foi 40,0 para todas as regiões, com exceção da Mediana na região Norte ( 0,0 ), o que indica que, nessa região, pelo menos metade dos estudantes tiraram nota zero nessa questão.
A nota Mínima foi zero para todas as regiões.
O Gráfico 6.12 representa a distribuição de notas da questão discursiva 4, do Componente de Conhecimento Específico.
Essa distribuição tem moda na classe de estudantes que deixaram a resposta dessa questão em branco nos cursos de Licenciatura e no intervalo ( 30 ; 40 ], para os de Bacharelado. 6.3.2.4 Comentários sobre as respostas à Questão Discursiva 4 A questão tratou do tema da aprovação pelo CONABIO das Metas Nacionais de Biodiversidade ( 2011-2020 ), sendo clara, direta e com pouca margem a interpretações diferentes da esperada.
A figura apresentada no enunciado da questão era de fácil entendimento, mostrando o percentual existente de área protegida nos diferentes biomas brasileiros, e a meta percentual a ser alcançada.
A linguagem utilizada foi bastante clara, e a questão tinha baixo nível de dificuldade quanto ao conteúdo exigido, sendo tema corrente na mídia.
O conteúdo abordado encontra -se no campo da Ecologia, embora, de modo tangente, mais claramente com enfoque na preservação e conservação de biomas, tema pertinente e importante para os graduandos em Ciências Biológicas ( Bacharelado e Licenciatura ).
Para uma questão pautada em valores de conservação da biodiversidade, o padrão de resposta foi bastante amplo.
O nível de profundidade exigido foi bastante básico, com muita flexibilidade, correspondendo a uma previsão de grande diversidade de respostas entre os estudantes.
No entanto, as respostas dos estudantes foram bem mais conservadoras do que as estabelecidas no padrão de respostas.
Não houve resposta correta que não constasse do padrão.
Nessa questão, dentre os graduandos que buscaram responder à questão, os de Licenciatura apresentaram respostas melhores do que os de Bacharelado.
Observa -se que a distribuição das notas de Licenciatura tem frequência crescente, enquanto para o Bacharelado houve uma concentração maior de notas até 40.
Como a questão era sobre o meio ambiente e os estudantes de Bacharelado são da área de Saúde, Biotecnologia e de Meio Ambiente, talvez apenas os últimos tenham respondido à questão de modo satisfatório.
Já na Licenciatura, a temática do Meio Ambiente é bastante recorrente, uma vez que se trata de tema de abrangência mundial e de fundamental importância para os que vão atuar na formação de jovens.
Os estudantes do Bacharelado e da Licenciatura tiveram desempenho melhor nessa questão, em comparação com a Questão 3.
Os estudantes, em geral, responderam bem à questão, apesar se ser comum e excessivo o uso de jargões de imprensa, fora da linguagem científica, algumas vezes se equivocando na abordagem dos conceitos.
No item ‘ a ’, as respostas mais comuns foram “ a proteção de espécies ameaçadas de extinção ”, “ recuperar ecossistemas degradados ”,"proporcionar meios e incentivos para a pesquisa científica, etc.".
No item ‘ b ’, a resposta mais frequente foi “ favorecer a participação popular e incentivar a alfabetização ( ou divulgação científica ) ”.
Apesar de um resultado melhor, a questão não permitiu diagnosticar as deficiências de conhecimentos dos estudantes, uma vez que a maioria das respostas foram superficiais, de conteúdo inespecífico.
Dessa forma, como a questão envolvia a interpretação de dados disponibilizados no enunciado e como houve um número elevado de respostas em branco, pode -se indicar, apenas, a necessidade de os cursos aprofundarem as experiências dos estudantes envolvendo a interpretação de tabelas, figuras, gráficos e, em particular, a proposição de experiências e projetos, estabelecendo objetivos, relevância, resultados e discussão acadêmica. 6.3.2.5 Análise da Questão Discursiva 5 do Componente de Conhecimento Específico Diferentemente das duas questões discursivas anteriores, a questão 5 da prova de Bacharelado era diferente da que foi apresentada na prova de Licenciatura.
A questão proposta aos bacharelandos envolvia conhecimentos sobre o estabelecimento dos fatores que atuam na zonação de populações marinhas.
Já a questão 5 apresentada na prova dos licenciandos envolvia conhecimentos na área de saúde e sua contribuição direta para as comunidades da escola e as fora dela.
As Tabelas 6.20a e 6.20b contém, respectivamente para Licenciatura e Bacharelado, as informações relativas à questão 5 do Componente de Conhecimento Específico.
O desempenho dos estudantes de ambas as Habilitações nessas questões foi o mais alto dentre as três questões discursivas desse componente.
As Médias para os alunos de Licenciatura ( ver Tabela 6.20a ) foram, respectivamente para os de cursos Presenciais e a Distância, 57,5 e 54,7 ( 57,0 para o total ).
A nota Máxima ( 100,0 ) foi atingida por, pelo menos, um concluinte de todas as regiões para ambas a Modalidades de Ensino.
A Mediana foi 70,0 para os cursos Presenciais em quase todas as regiões, com exceção da Mediana nas regiões Norte e Centro-Oeste ( 60,0 ).
Já para os cursos a Distância, a Mediana foi 60,0 em quase todas as regiões, com exceção da mediana nas regiões Norte ( 40,0 ) e Sudeste ( 70,0 ).
A nota Mínima foi zero para todas as combinações de Modalidade de Ensino e região.
A Média para os alunos cursos Presenciais de Bacharelado ( ver Tabela 6.20b ) foi 31,6.
Nas regiões Norte ( 22,0 ), Nordeste ( 30,2 ), Sul ( 31,2 ) e Centro-Oeste ( 28,5 ) a Média foi menor e na região Sudeste ( 33,4 ) maior que a do Brasil como um todo.
A nota Máxima ( 100,0 ), a Mediana ( 25,0 ) e nota Mínima ( 0,0 ) foram as mesmas para todas as regiões.
O Gráfico 6.13 representa a distribuição das notas da questão discursiva 5 do Componente de Conhecimento Específico.
Essa distribuição tem moda principal nas notas do intervalo ( 70 ; 80 ] para a Licenciatura e nas notas do intervalo ( 20 ; 30 ] para o Bacharelado. 6.3.2.6 Comentários sobre as respostas à Questão Discursiva 5 Questão 5 – Bacharelado A questão levou o aluno de Bacharelado em Ciências Biológicas a propor uma experiência que permitisse estabelecer os principais fatores abióticos e biológicos atuantes na zonação de populações de cracas ao longo de um costão rochoso.
O texto motivador era de fácil leitura, com linguagem bastante direta e clara.
O conteúdo abordado, segundo as Diretrizes Curriculares do MEC para a área de Ciências Biológicas, encontra -se em “ Ecologia : relações entre os seres vivos e destes com o ambiente ao longo do tempo geológico.
Conhecimento da dinâmica das populações, comunidades e ecossistemas, da conservação e manejo da fauna e flora e da relação saúde, educação e ambiente ”.
Portanto, o estabelecimento dos fatores que atuam na zonação de populações marinhas é bastante pertinente e importante para os graduandos do Bacharelado em Ciências Biológicas.
No entanto, como o tema, mais uma vez, voltou -se para o Meio Ambiente, pode ter sido de maior dificuldade para os alunos do Bacharelado em Saúde ou em Biotecnologia.
O nível de profundidade exigido estava acima do que se pode esperar dos graduandos.
Além disso, muitos não entenderam, corretamente, o que estava sendo pedido.
A questão foi de dificuldade mediana, pois exigia a elaboração de uma experiência para comprovação de resultados.
O padrão de resposta era bastante amplo e flexível para abarcar diferentes respostas.
A questão avaliava a capacidade do aluno de propor uma experiência e uma avaliação da importância dos fatores bióticos e abióticos na zonação desses invertebrados.
A necessidade de elaboração de uma experiência tornou a questão realmente interessante para os graduandos do Bacharelado, no entanto, a questão era mais simples para os mais envolvidos em trabalhos de pesquisa.
Os alunos do Bacharelado tiveram um desempenho mediano nessa questão, em comparação com as demais questões discursivas de Conhecimentos Específicos.
A maioria dos estudantes respondeu sobre os fatores abióticos e biológicos que atuam na zonaação, mas uma grande parte deles não propôs uma experiência comprovatória.
Tanto as respostas parciais quanto as melhores falharam na proposição da experiência.
Se o comando da questão tivesse feito menção explícita ao que se esperava de uma proposta de experiência, ou um projeto, com objetivo e resultados esperados, talvez houvesse um maior número de acertos.
As respostas com notas mais baixas foram aquelas que não propuseram uma experiência para o estabelecimento de quais fatores ambientais eram os mais importantes para as populações e, de uma maneira geral, não acertaram os fatores ambientais envolvidos na zonação das populações de cracas.
As respostas medianas foram aquelas que acertaram os fatores ambientais envolvidos, mas não propuseram uma experiência ou propuseram uma experiência fora do contexto.
As respostas boas foram aquelas que estabeleciam os fatores ambientais responsáveis pela zonação das cracas e propunham uma experiência viável no campo ou no laboratório.
Questão 5 – Licenciatura A Questão 5 para os cursos de Licenciatura foi bastante simples, com enunciado claro e de fácil entendimento para os estudantes.
O conteúdo abordado estava bastante adequado à Licenciatura, envolvendo conhecimento na área de saúde e sua contribuição direta para as comunidades da escola e fora dela.
A questão pode ser considerada fácil e não havia possibilidade de interpretações diferentes da esperada.
Os conhecimentos exigidos fazem parte das Diretrizes Curriculares da área Ciências Biológicas, assim como os conteúdos da Educação Básica, consideradas as Diretrizes Curriculares Nacionais para a formação de professores em nível superior, bem como as Diretrizes Nacionais para a Educação Básica e para o Ensino Médio.
O padrão de resposta estava bem claro, muito abrangente, com linguagem simples e bem adequada à questão.
O nível de exigência era básico, sendo, sem dúvida, a questão mais fácil e mais bem elaborada dentre as questões discursivas de Conhecimentos Específicos da prova.
No entanto, as respostas dos alunos foram muito menos diversificadas do que o esperado.
No item ‘ a ’, a resposta mais frequente foi a que apontava a capacidade que tem o professor de multiplicar as informações, de dar o exemplo de valores e de hábitos saudáveis e de ter o conhecimento correto sobre práticas de boa saúde.
No item ‘ b ’, os estudantes as respostas, de forma frequente, envolviam feiras de Ciências, sobre levar o conhecimento para a comunidade, fazer cartazes, dentre outros.
Os alunos responderam conforme o esperado pelo padrão de resposta e, de uma maneira geral, utilizaram uma linguagem adequada para a área.
Os alunos da Licenciatura tiveram um bom desempenho nessa questão, em comparação com as demais questões.
Os estudantes demonstraram reconhecerem o papel do professor como agente capaz de propagar conhecimentos e de mobilizar a escola e a comunidade para o debate de temas de saúde, de interesse de toda a coletividade.
Demonstraram em suas respostas conhecer bem o papel de liderança do professor na escola e na comunidade e souberam elaborar um evento para trabalhar junto com os alunos na prevenção de doenças.
A questão estava muito bem elaborada, o que \textbf{ajudou} muito no desempenho dos licenciandos. 6.3.3 Considerações Finais Os resultados dos graduandos de Bacharelado e de Licenciatura nessa prova revelaram alguns pontos importantes.
Na Questão 3, objetiva e clara, mas de conteúdo básico, os estudantes não foram muito bem.
Os conteúdos de Biologia Molecular, Genética e Bioquímica são estudados nos primeiros anos da graduação e, de acordo com o resultado, vão sendo esquecidos durante o curso.
Os estudantes de Bacharelado tiveram desempenho um pouco melhor nas suas respostas, talvez pelo fato do conteúdo ser mais repetido em outras disciplinas do curso durante a formação dos bacharéis.
Na Questão 4, mais subjetiva e com um padrão de resposta mais aberto, os graduandos tiveram um desempenho melhor.
No entanto, não se exigia conhecimentos científicos foi frequente o uso de uma linguagem pouco científica para a área.
Na correção, os estudantes de Licenciatura tiveram desempenho melhor, mas se restringiram ao conteúdo de divulgação, observado nas mídias.
Nas questões 5, mais relacionadas à área de atuação específica dos graduandos de Bacharelado e de Licenciatura, foi possível observar a relação direta que existe entre a forma como uma questão está estruturada e o resultado alcançado pelos estudantes.
Os estudantes de Bacharelado não entenderam bem o que havia sido pedido na Questão 5, e o resultado ficou muito aquém do esperado.
No entanto, na Questão 5 da Licenciatura, cujo enunciado era claro, \textbf{ajudando} no planejamento da resposta, o resultado foi muito mais satisfatório.

% matched lemmas: ajudar, localizar
\end{textsample}
